\documentclass[11pt,a4paper]{article}

\usepackage[utf8]{inputenc}
\usepackage[T1]{fontenc}
\usepackage{lmodern}
\usepackage{geometry}
\usepackage{setspace}
\usepackage{amsmath,amssymb,amsthm}
\usepackage{booktabs}
\usepackage{longtable}
\usepackage{pdflscape}
\usepackage{array}
\usepackage{enumitem}
\usepackage{natbib}
\setcitestyle{aysep={}}
\usepackage{hyperref}
\usepackage{xcolor}
\usepackage{caption}
\usepackage{float}
\usepackage{authblk}
\usepackage{url}

\geometry{margin=1in}
\doublespacing
\hypersetup{
    colorlinks=true,
    linkcolor=blue!50!black,
    citecolor=blue!50!black,
    urlcolor=blue!50!black
}
\setlist[itemize]{leftmargin=*}
\setlist[enumerate]{leftmargin=*}
\setlength{\parskip}{0.35em}

\newtheorem{definition}{Definition}
\newtheorem{proposition}{Proposition}
\newtheorem{assumption}{Assumption}
\newtheorem{lemma}{Lemma}

\newcommand{\GDP}{Y}
\newcommand{\FS}{FS}
\newcommand{\MFC}{MFC}
\newcommand{\LS}{LS}
\newcommand{\AI}{AI}
\newcommand{\Prob}{\Pr}

\title{Online Appendix for: AI, Fiscal Space, and Universal Social Protection in Emerging Economies: A Robust Accounting-Fiscal Threshold Framework}
\author{Vicenzo Scavino}
\affil{}
\date{}

\usepackage{xr-hyper}
\externaldocument[][nocite]{main}

\begin{document}
\maketitle

\noindent This supplement accompanies the main article; numbered cross-references point to the main text.

\appendix
\section{AI-Shock Construction Module}\label{app:ai_shock_module}

This appendix formalizes the robustness checks behind the AI-shock block. The baseline text uses the frontier-normalized multiplicative translation factor because it captures complementarity and avoids dependence on internal-sample statistics. None of the auxiliary choices below is treated as structural.

\subsection{Normalization of input variables}

The preparedness reconstruction from the AIPI's four dimensions is
\begin{equation}
    A_{i,t}=\left(D_{i,t}^{w_D}H_{i,t}^{w_H}R_{i,t}^{w_R}G_{i,t}^{w_G}\right)^{1/(w_D+w_H+w_R+w_G)},
\end{equation}
Here $D$ is digital infrastructure, $H$ human capital and labour-market policy readiness, $R$ innovation and economic integration, and $G$ regulation and ethics. Geometric aggregation makes a severe bottleneck reduce absorptive capacity even when another dimension is strong; reconstructed and published AIPI values need not coincide. These symbols are local: elsewhere $D_{s,t}$ is diffusion, $H$ an accumulation horizon, and $R$ denotes revenue objects.

All AI-shock inputs must be transformed to a common $[0,1]$ scale before aggregation. The simplest transformation is min--max normalization:
\begin{equation}
    X^{MM}_{i,t}=\frac{X_{i,t}-\min_i X_{i,t}}{\max_i X_{i,t}-\min_i X_{i,t}}.
\end{equation}
This expression is used only when $\max_iX_{i,t}>\min_iX_{i,t}$. If all observations are equal, min--max normalization is not informative; a diagnostic value such as $0.5$ may be assigned only for descriptive plots, while threshold calculations must use external-frontier normalization. Because min--max normalization can depend strongly on sample composition and outliers, percentile-rank and frontier normalizations are available as labelled robustness alternatives:
\begin{align}
    X^{pct}_{i,t}&=\frac{rank(X_{i,t})-1}{N_t-1}\quad (N_t>1),\\
    X^{frontier}_{i,t}&=\min\left\{1,\frac{X_{i,t}}{X^{frontier}_t}\right\}.
\end{align}
The frontier reference is preferred to the maximum because it is less sensitive to a single outlier when constructed from an external benchmark set. For the baseline AI-translation factor, $X^{frontier}_t$ must be external to the pilot emerging-economy sample and should represent a high-preparedness, high-adoption, high-exposure benchmark. An internal pilot-sample percentile is a relative-position diagnostic, not an absolute translation benchmark. Single-country pilots must not use min--max or percentile-rank normalizations for threshold construction; they must use external-frontier normalization. Share-valued inputs that enter the translation score directly, $E^{prod}$ and $q^{prod}$, are already measured on the natural $[0,1]$ scale and enter $S$ in natural units: the only frontier normalization in the baseline translation factor is the score-level division by $S^{frontier}_{s,t}$ in Equation \eqref{eq:T_ndc_main_v10}. The input-level normalizations of this subsection apply to index-type inputs such as the preparedness dimensions and $Gap$. Applying an input-level frontier normalization to $E^{prod}$ or $q^{prod}$ and then dividing the score by $S^{frontier}$ would discount the same distance to the frontier twice and is not admissible.

\subsection{Exposure construction from occupational and sectoral maps}

When both maps are available, productive exposure is constructed as
\begin{align}
    E^{occ,prod}_{i,t}&=\sum_o \omega^{emp}_{i,o,t}e^{prod}_o,\\
    E^{sec,prod}_{i,t}&=\sum_j \omega^{va}_{i,j,t}e^{prod}_j,
\end{align}
and combined by
\begin{equation}
    E^{prod}_{i,t}=\varpi E^{occ,prod}_{i,t}+(1-\varpi)E^{sec,prod}_{i,t},
    \qquad \varpi\in[0,1].
\end{equation}
Employment weights the occupational map and value added the sectoral map; $\varpi$ selects their convex combination.

\subsection{Technical floor}

The regularized input is
\begin{equation}
    \tilde X_{i,t}=\epsilon+(1-\epsilon)X_{i,t},
\end{equation}
with
\begin{equation}
    \epsilon\in\{0,0.01,0.05\}.
\end{equation}
For low but positive interior adoption, the floor also creates a level distortion: with $\epsilon=0.05$, a small $q^{prod}$ is inflated proportionally by the transformation. The grid therefore includes $\epsilon=0$, which is the preferred option whenever threshold inversion is numerically stable. The baseline implementation fixes the numerical floor at $\epsilon=10^{-9}$ for support checks and de-flooring, the numerically stabilized version of the preferred grid point $\epsilon=0$, with the logit-inversion floor fixed at $\epsilon_\mu=10^{-6}$.
The floor is a numerical regularization device, not a substantive claim that every country has positive AI capability in all dimensions. The implementation therefore separates the numerical object $\tilde q^{prod,num}=\epsilon+(1-\epsilon)q^{prod}$ from the substantive adoption object $q^{prod}$ and the substantive exposure object $E^{prod}$. The floor is never allowed to create translation or fiscal space by itself. In every translation score, multiplicative or additive, zero substantive adoption or zero substantive productive exposure is imposed as a hard gate:
\begin{equation}\label{eq:q_zero_gate_v21}
    q^{prod}_{i,s,t}\leq0\ \text{or}\ E^{prod}_{i,t}\leq0\quad\Longrightarrow\quad
    S_{i,s,t}=0,\quad T_{i,s,t}=0,\quad g^{AI}_{i,s,t}=0.
\end{equation}
Equivalently, the multiplicative score is written with the indicator $\mathbf{1}\{q^{prod}>0,E^{prod}>0\}$. The regularized values $\tilde q^{prod,num}$ and $\tilde E^{prod}$ may stabilize elasticities or threshold inversions conditional on positive substantive adoption and exposure, but they cannot convert zero adoption or zero exposure into positive AI-induced output. The gate removes crossings with substantive adoption exactly equal to zero, but floor-assisted crossings remain possible when $q^{prod}$ is positive but very small. Two complementary checks therefore remain active in the baseline implementation: the $FloorDriven$ diagnostic, under its unique operational definition in Section \ref{subsec:impossibility_score}, Equation \eqref{eq:floor_driven_operational_v25}, flags truncated cells with $q^{prod,*}=0$ whose crossing depends on $\epsilon>0$, while floor sensitivity of crossings with small positive adoption is diagnosed by re-running the threshold inversions at the preferred grid point $\epsilon=0$.


\subsection{Alternative aggregation forms}

The baseline no-double-counting multiplicative index is
\begin{equation}
    S^{NDC}_{i,s,t}=\mathbf{1}\{q^{prod}_{i,s,t}>0,\ E^{prod}_{i,t}>0\}(\tilde E^{prod}_{i,t})^{\beta}(\tilde q^{prod,num}_{i,s,t})^{\gamma}.
\end{equation}
This is the preferred form when productive adoption is simulated from preparedness and frictions, $q^{prod}=h(A,I,Gap)$. Preparedness affects the shock through adoption and therefore is not multiplied again in the baseline translation index.

A full-score robustness specification is
\begin{equation}
    S^{full}_{i,s,t}=\mathbf{1}\{q^{prod}_{i,s,t}>0,\ E^{prod}_{i,t}>0\}\tilde A_{i,t}^{\alpha}(\tilde E^{prod}_{i,t})^{\beta}(\tilde q^{prod,num}_{i,s,t})^{\gamma}.
\end{equation}
This form is valid only when $q^{prod}$ is observed independently from adoption data, or when it is explicitly labeled as a full-score robustness check rather than a no-double-counting baseline. If the adoption equation includes $E^{prod}$ directly, the direct exposure term must be removed from the baseline score and the score becomes
\begin{equation}
    S^{q-only}_{i,s,t}=\mathbf{1}\{q^{prod}_{i,s,t}>0,\ E^{prod}_{i,t}>0\}(\tilde q^{prod,num}_{i,s,t})^{\gamma}.
\end{equation}

An additive robustness specification is
\begin{equation}
    S^{add}_{i,s,t}=\mathbf{1}\{q^{prod}_{i,s,t}>0,\ E^{prod}_{i,t}>0\}\left(w_E E^{prod}_{i,t}+w_q q^{prod}_{i,s,t}\right),
    \qquad w_E+w_q=1,
\end{equation}
with preparedness added only when adoption is independently observed. The zero-adoption/zero-exposure gate applies to the additive form exactly as to the multiplicative forms: without the indicator, a robustness run could report positive translation and fiscal space with zero substantive adoption, which Equation \eqref{eq:q_zero_gate_v21} rules out by construction. The additive form is not the preferred baseline because, conditional on passing the gate, it still permits one strong component to compensate for a nearly absent component.

\subsection{Elasticity grid}

Since the structural elasticities governing the translation of AI exposure into aggregate productivity are not directly observed, the baseline uses unit elasticities for transparency. The named specifications below span the main directions; the full sensitivity grid for active elasticities is $\{0.5,1,1.5,2\}$, as declared in the master calibration matrix. The elasticity menu is:
\begin{table}[H]
\centering
\caption{Elasticity-grid robustness for AI-shock translation}
\label{tab:elasticity_grid}
\small
\begin{tabular}{lccc p{0.34\textwidth}}
\toprule
Specification & $\alpha$ & $\beta$ & $\gamma$ & Interpretation \\
\midrule
No-double-counting baseline & 0 & 1 & 1 & Preparedness affects adoption, not the score directly \\
Exposure-heavy & 0 & 1.5 & 1 & Occupational/sectoral opportunity matters more \\
Adoption-heavy & 0 & 1 & 1.5 & Realized diffusion matters more \\
Low complementarity & 0 & 0.5 & 0.5 & Lower penalty for weak components \\
High complementarity & 0 & 1.5 & 1.5 & Stronger penalty for weak components \\
Full-score robustness & 1 & 1 & 1 & Only when adoption is independently observed \\
\bottomrule
\end{tabular}
\end{table}

\subsection{No-double-counting checks}

The implementation must report which information enters each component. The recommended checks are:
\begin{align}
    S^{NDC}_{i,s,t}&=\mathbf{1}\{q^{prod}_{i,s,t}>0,\ E^{prod}_{i,t}>0\}(\tilde E^{prod}_{i,t})^{\beta}(\tilde q^{prod,num}_{i,s,t})^{\gamma},\qquad q^{prod}=h(A,I,Gap),\\
    S^{q-only}_{i,s,t}&=\mathbf{1}\{q^{prod}_{i,s,t}>0,\ E^{prod}_{i,t}>0\}(\tilde q^{prod,num}_{i,s,t})^{\gamma},\qquad q^{prod}=h(A,E^{prod},I,Gap),\\
    S^{full}_{i,s,t}&=\mathbf{1}\{q^{prod}_{i,s,t}>0,\ E^{prod}_{i,t}>0\}\tilde A_{i,t}^{\alpha}(\tilde E^{prod}_{i,t})^{\beta}(\tilde q^{prod,num}_{i,s,t})^{\gamma},\qquad q^{prod}\ \text{observed independently.}
\end{align}
The first is the baseline. The second is used if exposure enters adoption. The third is a robustness specification, not the no-double-counting baseline.

\subsection{Scenario timing as robustness only}

The baseline paper is not designed to estimate calendar dates for threshold crossing. If a dynamic diffusion path is used, it should be treated only as a robustness extension:
\begin{equation}
    g^{AI}_{i,s,t}=\phi^{Y,nom}_s D_{s,t}T_{i,s,t},
    \qquad g^{AI,max}_{s,t}=\phi^{Y,nom,max}_sD_{s,t}T^{max},\quad T^{max}\equiv1,
\end{equation}
where
\begin{equation}
    D_{s,t}=\frac{1}{1+\exp[-k_s(t-t^0_s)]},
\end{equation}
and $D_{s,t}=1$ in the static implementation.
This extension changes the annual path of the shock but does not alter the central accounting-fiscal threshold logic.

\subsection{Robustness menu}

The framework's robustness menu spans the following dimensions; the executed subset and its outputs are recorded in the prespecified plans and the robustness summary of the replication package:
\begin{longtable}{p{0.25\textwidth}p{0.62\textwidth}}
\caption{AI-shock robustness dimensions}\label{tab:ai_shock_robustness}\\
\toprule
Dimension & Variants \\
\midrule
\endfirsthead
\toprule
Dimension & Variants \\
\midrule
\endhead
Normalization & Min--max, percentile rank, external frontier; internal-sample percentile only as relative-index robustness \\
Functional form & Multiplicative baseline, additive robustness \\
Elasticities & No-double-counting baseline, exposure-heavy, adoption-heavy, low-complementarity, high-complementarity, full-score only when adoption is independently observed \\
Base scenario $\phi^{Y,nom}_s$ & Conservative, moderate, high, disruptive stress test \\
Adoption & Observed adoption, externally calibrated adoption, simulated logistic adoption \\
Exposure & Occupational, sectoral, task-based, productive versus labor-displacement exposure \\
Double counting & $S^{NDC}$ baseline, $S^{q\text{-}only}$ when exposure enters adoption, $S^{full}$ robustness \\
Temporal structure & Static scenario, logistic diffusion robustness only \\
\bottomrule
\end{longtable}

The methodological claim is conditional: given an externally disciplined AI-productivity scenario, the model evaluates how much of that scenario can plausibly translate into domestic fiscal space under alternative assumptions about exposure, preparedness, adoption, and fiscal capture.




\section{Productive AI Adoption Module}\label{app:adoption_module}

This appendix gives the reproducible adoption module. The baseline separates exposure from adoption: exposure determines where AI can affect tasks; adoption determines whether firms and workers actually integrate AI into production.

\subsection{Access, use, and productive adoption}

Let $q^{access}$ denote access, $q^{use}$ denote use, and $q^{prod}$ denote productive adoption. The ordering is
\begin{equation}
    0\leq q^{prod}_{i,s,t}\leq \bar q_{i,t}\leq1,
\end{equation}
with $q^{access}$ and $q^{use}$ reported when observed. The stronger inequality $q^{prod}\leq q^{use}$ is imposed only if the empirical design treats both as same-period flow variables or assumes non-decreasing use. In the baseline, $q^{prod}$ is stock-like organizational integration, so it may exceed current use after a temporary decline in use.

\subsection{Logistic diffusion equation}

The baseline logistic use equation is
\begin{equation}
    q^{use}_{i,s,t}=\bar q_{i,t}\Lambda(\ell^q_{i,s,t}),
    \qquad
    \ell^q_{i,s,t}=\mu_{i,s,t}+\nu_A A_{i,t}-\nu_I I_{i,t}-\nu_G Gap_{i,t}.
\end{equation}
Productive exposure $E^{prod}$ is deliberately excluded from this baseline adoption equation because it already enters the translation score $S^{NDC}$. A robustness version may include $\nu_EE^{prod}$ in $\ell^q$, but then the direct $E^{prod}$ term must be removed from $S$.

\subsection{Country-specific adoption ceiling}

The ceiling is
\begin{equation}
    \bar q_{i,t}=\max\{0,\min[1,1-\omega_I I_{i,t}-\omega_G Gap_{i,t}]\}.
\end{equation}
For the economic rationale behind the informality term, see Section \ref{subsec:translation_factor_main} in the main text. A convex ceiling variant, for example $\bar q=(1-I)^{a_I}(1-Gap)^{a_G}$, may be reported as a functional-form robustness check; the baseline linear ceiling is a convention, not an estimated form.
The ceiling uses observed or harmonized informality and digital/skills gaps. Informality and $Gap$ may affect both diffusion speed and long-run capacity only if their empirical content is partitioned or residualized; otherwise the same barrier is counted twice. If a target adoption moment falls outside the interior support, the target is adjusted:
\begin{equation}
    q^{use,target,adj}_{i,t_0}
    =\min\left\{\max\left\{q^{use,target}_{i,t_0},\epsilon_\mu\right\},\bar q_{i,t_0}-\epsilon_\mu\right\},
    \qquad \bar q_{i,t_0}>2\epsilon_\mu.
\end{equation}
If $\bar q_{i,t_0}\leq2\epsilon_\mu$, the target moment is not logit-invertible and the cell is labelled ceiling-impossible for that calibration.

\subsection{Initial adoption moment}

The initial intercept is
\begin{equation}
    \mu_{i,s,t_0}=\log\left(\frac{q^{use,target,adj}_{i,t_0}}{\bar q_{i,t_0}-q^{use,target,adj}_{i,t_0}}\right)
    -\nu_A A_{i,t_0}+\nu_I I_{i,t_0}+\nu_G Gap_{i,t_0}.
\end{equation}
The calibration registry records the source of $q^{use,target}_{i,t_0}$: an OECD--Eurostat benchmark adoption average scaled by country AIPI relative to the benchmark AIPI, with a declared $\pm30\%$ support. This is an anchored diffusion convention: $\mu_{i,s,t_0}$ is a residual initial-condition term, not evidence that $A$, $I$, and $Gap$ have been structurally identified as predictors of the initial level. If no adoption data are available, low/central/high targets are calibrated from digital access, firm digitalization, and country preparedness. If the paper wants a predictive adoption model, it must use a separate structural-intercept robustness specification with $\mu_s$ or $\mu_s+\alpha_i$ and evaluate fit without target inversion. The module contains no endogenous imitation or network-diffusion term of the Bass type: with the anchored intercept held fixed, adoption paths are driven by scenario updates to preparedness, frictions, gaps, and ceilings. Adoption trajectories are therefore scenario paths, not endogenous-diffusion forecasts; if GenAI diffuses through imitation faster than preparedness improves, the baseline understates medium-term adoption.

\subsection{Productive-adoption lag}

The productive adoption path is
\begin{equation}
    q^{prod}_{i,s,t}=\min\left\{\bar q_{i,t},(1-\lambda_q)q^{prod}_{i,s,t-1}+\lambda_q q^{use}_{i,s,t}\right\},
    \qquad \lambda_q\in\{0.15,0.30,0.50\}\quad\text{in baseline, with }\lambda_q=1\text{ only as an instant-adoption stress test}.
\end{equation}
The default initial condition is $q^{prod}_{i,s,t_0}=q^{use,target,adj}_{i,t_0}$, with $q^{prod}_{i,s,t_0}=0$ used as a conservative robustness check. The value $\lambda_q=1$ is an immediate-pass-through stress test, not the baseline.

\subsection{Coefficient restrictions and robustness}

The baseline restrictions are
\begin{equation}
    \nu_A\geq0,\qquad \nu_I\geq0,\qquad \nu_G\geq0,
\end{equation}
where the signs are written so that informality and gaps lower adoption through minus signs in $\ell^q$. If the robustness equation includes $\nu_EE^{prod}$, then $\nu_E\geq0$ and the direct exposure term is removed from $S^{NDC}$ to preserve the no-double-counting convention.

\section{Functional Income Allocation Module}\label{app:labor_share_module}

This appendix formalizes the robustness checks behind the labor-share channel in Section \ref{sec:labor_share}. The module is not a wage-forecasting model. It is an accounting device for allocating AI-induced output between labor and capital/profit channels so that the fiscal conversion block can apply different taxability assumptions to each channel.

\subsection{Raw and adjusted labor share}

The preferred empirical input is an adjusted labor share. In economies with substantial self-employment or informal work, the observed compensation-of-employees share can understate the labor component of income because mixed income combines labor and capital returns. A defensible implementation should report:
\begin{equation}
    LS^{raw}_{i,t}=\frac{\text{Compensation of employees}_{i,t}}{Y^0_{i,t}},
\end{equation}
plus an adjusted measure $LS^{adj}_{i,t}$ when mixed-income imputations are available. The main results state the version used: the country-specific PWT labor share, recorded in the calibration registry; ranking stability under raw versus adjusted measures was tested in a labelled post-baseline robustness extension using the raw UN SNA compensation-of-employees share as the alternative input, with effective rates re-derived consistently: policy rankings and country-policy classifications are unchanged across the full grid, and the only cell-level effect is that Peru's most marginal stress-scenario pension crossing, in the status quo regime, falls just below the accounting threshold ($V^{gross}$ from $1.02$ to $0.98$), quantifying the fragility already flagged (amendment registry, replication package).

\subsection{Automation, augmentation, and reskilling scenarios}

The latent factor-share equation separates three mechanisms:
\begin{equation}
    \Delta \ell_{i,s,t}
    =
    -\left(\delta_s-\zeta_s RST_{i,t}\right)E^{auto}_{i,t}q^{prod}_{i,s,t}
    +\psi_s E^{aug}_{i,t}q^{prod}_{i,s,t}.
\end{equation}
The baseline restrictions are
\begin{equation}
    \delta_s,\psi_s,\zeta_s\geq0,
    \qquad
    0\leq \lambda_{LS}\leq1.
\end{equation}
The empirical implementation should treat these as scenario parameters, not as known structural constants.

\begin{longtable}{p{0.25\textwidth}p{0.11\textwidth}p{0.11\textwidth}p{0.14\textwidth}p{0.27\textwidth}}
\caption{Labor-share scenario logic}\label{tab:labor_share_scenarios}\\
\toprule
Scenario & $\delta_s$ & $\psi_s$ & $\zeta_s$ & Interpretation \\
\midrule
\endfirsthead
\toprule
Scenario & $\delta_s$ & $\psi_s$ & $\zeta_s$ & Interpretation \\
\midrule
\endhead
Automation-heavy & High & Low & Low & AI mainly substitutes tasks \\
Balanced & Medium & Medium & Medium & Substitution and complementarity coexist \\
Augmentation-heavy & Low & High & Medium/high & AI mainly complements labor \\
High-reskilling & Medium & Medium & High & Institutions absorb displacement pressure \\
\bottomrule
\end{longtable}


\subsection{Exact identities and first-order approximation}

The exact labor-income and non-labor-residual changes are
\begin{equation}
    \frac{\Delta W^{AI}}{Y^0}=LS'(1+g^{AI})-LS,
\end{equation}
\begin{equation}
    \frac{\Delta NLS^{AI}}{Y^0}=NLS'(1+g^{AI})-NLS,
    \qquad NLS=1-LS,\quad NLS'=1-LS'.
\end{equation}
Only a share $\chi^{Kbase}$ of $\Delta NLS^{AI}$ is treated as taxable capital/profit base:
\begin{equation}
    \frac{\Delta \Pi^{AI,gross}}{Y^0}=\chi^{Kbase}\frac{\Delta NLS^{AI}}{Y^0},\qquad 0\leq\chi^{Kbase}\leq1.
\end{equation}
Since $\Delta LS^{AI}=LS'-LS$, expanding the labor equation gives
\begin{equation}
    \frac{\Delta W^{AI}}{Y^0}
    =LS\,g^{AI}+\Delta LS^{AI}+\Delta LS^{AI}g^{AI}.
\end{equation}
The first-order approximation is therefore
\begin{equation}
    \frac{\Delta W^{AI}}{Y^0}\approx LS\,g^{AI}+\Delta LS^{AI},
\end{equation}
when the interaction term $\Delta LS^{AI}g^{AI}$ is small. For the non-labor residual,
\begin{equation}
    \frac{\Delta NLS^{AI}}{Y^0}
    =NLS\,g^{AI}-\Delta LS^{AI}-\Delta LS^{AI}g^{AI},
\end{equation}
with first-order approximation
\begin{equation}
    \frac{\Delta NLS^{AI}}{Y^0}\approx NLS\,g^{AI}-\Delta LS^{AI}.
\end{equation}
The exact identities are the benchmark; the approximations are only interpretive shortcuts.

The parameter $\chi^{Kbase}$ converts the broad non-labor residual into the portion that is plausibly taxable as corporate profits, business income, capital income, or AI-related rents. Because it multiplies the AI-induced increment, $\chi^{Kbase}$ is a marginal-composition parameter: the AI increment is plausibly more concentrated in corporate and formal sectors than the economy-wide non-labor residual, so the taxable share of the increment may exceed the average taxable share of $NLS$, and calibrating $\chi^{Kbase}$ from average national-accounts composition is a conservative convention. Supports should therefore be chosen for the incremental base, mirroring the marginal-informality convention of Appendix \ref{app:informality_module}, with sectorally weighted variants reported where data allow. The labor and non-labor residual components exhaust the AI-induced output increment:
\begin{equation}
\begin{split}
    &\Bigl[LS'_{i,s,t}(1+g^{AI}_{i,s,t})-LS_{i,t}\Bigr]
    +\Bigl[NLS'_{i,s,t}(1+g^{AI}_{i,s,t})-NLS_{i,t}\Bigr] \\
    &=g^{AI}_{i,s,t}.
\end{split}
\end{equation}

\subsection{Bounds and functional-form robustness}

The baseline uses a logit transformation to keep $LS'$ in $(0,1)$. Define the share-point shock, measured in labor-share points, implied by the logit model as
\begin{equation}
    \Delta LS^{logit\text{-}eq}_{i,s,t}
    =\Lambda\left(\log\frac{LS_{i,t}}{1-LS_{i,t}}+\Delta\ell_{i,s,t}\right)-LS_{i,t}.
\end{equation}
Since the first term in the difference is $LS^{target}_{i,s,t}$, direct addition of this shock is algebraically identical to the baseline partial-adjustment rule:
\begin{equation}
    LS_{i,t}+\lambda_{LS}\Delta LS^{logit\text{-}eq}_{i,s,t}
    =LS_{i,t}+\lambda_{LS}(LS^{target}_{i,s,t}-LS_{i,t})
    =(1-\lambda_{LS})LS_{i,t}+\lambda_{LS}LS^{target}_{i,s,t}
    =LS'_{i,s,t}.
\end{equation}
Because $LS_{i,t},LS^{target}_{i,s,t}\in(0,1)$ and $\lambda_{LS}\in[0,1]$, this convex combination remains in $(0,1)$, so outer truncation never activates and is not an independent robustness check. The only functional-form robustness is a linear-share specification with separately calibrated linear-scale parameters $\delta^{lin}_s$, $\psi^{lin}_s$, and $\zeta^{lin}_s$; unlike the logit baseline, the linear specification is not self-bounding, so $LS'$ must there be explicitly truncated to $(0,1)$ and any activated truncation must be reported. Coefficients calibrated for log-odds space must not be reused as direct percentage-point labor-share shocks.

\subsection{Taxable labor-income channel under informality}

The labor-share module produces an income-flow change. The fiscal conversion block then determines how much of that flow is taxable. The following base-side representation is an alternative to the baseline rate-side convention in Section \ref{subsec:gross_effective_fiscal_conversion}; it is valid only if $\tau^{L,fisc,eff}$ does not apply the labor-informality penalty again. In high-informality economies, taxable labor income may be approximated by
\begin{equation}\label{eq:labor_informality_base_side_v25}
    \frac{\Delta W^{taxable}_{i,s,t}}{Y^0_{i,t}}
    =
    (1-I^{L,marg}_{i,s,t})^{\rho_L}
    \frac{\Delta W^{AI}_{i,s,t}}{Y^0_{i,t}},
\end{equation}
where $I^{L,marg}_{i,s,t}$ is marginal labor informality as defined in Section \ref{subsec:marginal_informality}; aggregate $I^L_{i,t}$ is used only when sectoral AI-shock weights are unavailable, and $\rho_L\geq0$ controls the strength of the penalty. This alternative must follow the assignment rule in Section \ref{subsec:gross_effective_fiscal_conversion} and Table \ref{tab:no_double_counting_informality}. A more detailed base-side implementation may split formal and informal labor income:
\begin{equation}
    \Delta W^{taxable}=\Delta W^{formal}+\omega^{inf}\Delta W^{informal},
    \qquad 0\leq \omega^{inf}<1.
\end{equation}
Here $\omega^{inf}$ is the taxable fraction of the informal labor-income increment, and the split is exhaustive: $\Delta W^{formal}+\Delta W^{informal}=\Delta W^{AI}$. This is reported as a fiscal-capture adjustment, not as a claim about the welfare effect of informality.

\subsection{Labor-share robustness table}

\begin{longtable}{p{0.24\textwidth}p{0.30\textwidth}p{0.34\textwidth}}
\caption{Labor-share robustness dimensions}\label{tab:labor_share_robustness}\\
\toprule
Block & Baseline & Robustness \\
\midrule
\endfirsthead
\toprule
Block & Baseline & Robustness \\
\midrule
\endhead
Labor-share measure & Adjusted $LS$ when available & Raw $LS$; alternative mixed-income imputations \\
Exposure & $E^{auto}$ and $E^{aug}$ separated & Single labor exposure as fallback \\
Functional form & Logit-bounded $LS^{target}$ with partial adjustment & Linear-share specification with separately calibrated $\delta^{lin}_s$, $\psi^{lin}_s$, and $\zeta^{lin}_s$ \\
Adjustment speed & Partial adjustment $\lambda_{LS}$ from the declared grid & Grid endpoints: immediate adjustment ($\lambda_{LS}=1$); slower adjustment (low $\lambda_{LS}$) \\
Scenarios & Balanced automation/augmentation & Automation-heavy; augmentation-heavy; high-reskilling \\
Parameters & Scenario grid for $\delta,\psi,\zeta$ & Monte Carlo draws from bounded distributions \\
Informality & Effective labor-tax penalty & Formal/informal labor split \\
\bottomrule
\end{longtable}

The core reporting rule is to show both distributional and income-flow objects: $\Delta LS^{AI}$, $\Delta W^{AI}/Y^0$, $\Delta NLS^{AI}/Y^0$, and $\Delta \Pi^{AI,gross}/Y^0$. This avoids the mistake of interpreting a lower labor share as automatically lower labor income in levels.


\section{Fiscal Conversion Module}\label{app:fiscal_conversion_module}

This appendix formalizes the fiscal conversion block in Section \ref{sec:fiscal_conversion}. The central convention is simple: effective tax rates already contain capturability adjustments. Capturability shares $\chi^j$ are used only to construct $\tau^{j,eff}$ from statutory or benchmark rates; they are not multiplied again in the main $MFC$ equation. Throughout this appendix, $\tau^{j,eff}$ abbreviates the fiscal-capture rates of the main text, $\tau^{L,fisc,eff}$, $\tau^{K,fisc,eff}$, $\tau^{C,eff}$, and $\tau^{AI,eff}$; disposable-income rates are always written with the $disp$ superscript.

\subsection{Nested specifications and no double counting}

The framework defines three nested fiscal-conversion specifications; the channel specification is the operative baseline mode:
\begin{longtable}{p{0.25\textwidth}p{0.39\textwidth}p{0.24\textwidth}}
\caption{Nested fiscal conversion specifications}\label{tab:fiscal_nested_specs}\\
\toprule
Specification & Formula & Use \\
\midrule
\endfirsthead
\toprule
Specification & Formula & Use \\
\midrule
\endhead
Average conversion & $fs^{avg}=\theta g^{AI}$ & Transparent benchmark \\
Buoyancy conversion & $fs^{buoy}=b\theta g^{AI}$ & Reduced-form macro benchmark \\
Channel gross conversion & $fs^{gross}=MFC^{gross}g^{AI}$ & Preferred gross revenue object \\
Effective fiscal space & $fs^{eff}=\widetilde{MFC}^{gross}g^{AI}-ac^{net,fix}-tr^{ann}-leak^{fix}$ & Feasibility benchmark \\
\bottomrule
\end{longtable}
The average and buoyancy formulas are never summed with the channel-based formula. They are alternative benchmarks for checking whether the preferred channel specification produces plausible magnitudes. Here $\theta$ denotes one of the observed revenue anchors $\theta^{tax}$, $\theta^{rev}$, or $\theta^{nonres}$ defined in Section \ref{app:relocated_main} of this supplement, and $b$ a historical buoyancy estimate $Buoyancy^{hist}$ from Section \ref{sec:historical_plausibility}; the anchor used must be stated.

\subsection{Statutory, average, effective, and marginal rates}

The relevant tax object for the AI shock is the effective marginal fiscal capture rate applied to the incremental tax base:
\begin{equation}
    \tau^{j,eff}_{i,s,r,t},\qquad j\in\{L,K,C,AI\}.
\end{equation}
When an observed effective marginal rate is available, it enters directly. When only a statutory or benchmark rate is available, the effective rate is constructed as
\begin{equation}\label{eq:tau_eff_decomposition_v10}
    \tau^{j,eff}_{i,s,r,t}=\tau^{j,base}_{i,r,t}\chi^j_{i,s,r,t}.
\end{equation}
After this construction, the main channel formula uses $\tau^{j,eff}$ only:
\begin{equation}
    MFC^{gross}=\sum_{j\in\{L,K,C\}}\omega^j\tau^{j,eff}+\lambda^{AI}\tau^{AI,eff}.
\end{equation}
It does not use $\tau^{j,eff}\chi^j$. Under the subset convention of Equation \eqref{eq:omega_k_net_v16}, the capital weight in this sum is $\omega^{K,net}$ rather than $\omega^K$.

\subsection{Channel-revenue function and frozen-weight convention}

\begin{equation}\label{eq:channel_revenue_function_v24}
\begin{split}
    rev^{ch}_{i,s,r,t}(g)
    &=\sum_{j\in\{L,K,C\}}\tau^{j,fisc,eff}_{i,s,r,t}\frac{\Delta B^{j,tax}_{i,s,t}(g)}{Y^0_{i,t}}
      +\lambda^{AI}_{i,s,t}(g)\tau^{AI,eff}_{i,s,r,t}g,\\
    fs^{eff,ch}_{i,p,s,r,t}(g)
    &=rev^{ch}_{i,s,r,t}(g)
      -g\left(ac^{ME}_{i,s,r,t}+tr^{ME}_{i,s,r,t}+leak^{ME}_{i,s,r,t}\right)
      -ac^{net,fix}_{i,p,s,r,t}-tr^{ann}_{i,s,r,t}-leak^{fix}_{i,s,r,t}.
\end{split}
\end{equation}
The reduced-form and frozen-weight versions are the special cases in which $rev^{ch}(g)$ is linear in $g$. Under the subset convention of Equation \eqref{eq:omega_k_net_v16}, $\Delta B^{K,tax}_{i,s,t}(g)$ in Equation \eqref{eq:channel_revenue_function_v24} is the domestic post-shifting profit base net of the AI-rent component, mirroring $\omega^{K,net}$; otherwise the rent flow would enter both the $K$ term and the AI-rent term. This distinction protects the threshold formulas below from treating an endogenous fiscal-conversion object as if it were fixed.

\subsection{Consumption channel robustness}

The baseline consumption channel is induced by income flows:
\begin{equation}
\begin{split}
    \Delta W^{AI,disp}&=(1-\tau^{L,disp})\Delta W^{AI},\\
    \Delta\Pi^{AI,disp}&=(1-\tau^{K,disp})\Delta\Pi^{AI,dom},\\
    \Delta R^{AI,gross}&=(1-\chi^{Kbase})\Delta NLS^{AI},\\
    \Delta R^{AI,disp}&=\theta^{R,dom}(1-\tau^{R,disp})\Delta R^{AI,gross}.
\end{split}
\end{equation}
\begin{equation}
    \Delta C^{AI,taxable}=(1-exempt^C)(1-m^M)\left(mpc^W\Delta W^{AI,disp}+mpc^\Pi\Delta\Pi^{AI,disp}+mpc^R\Delta R^{AI,disp}\right).
\end{equation}
The term $\Delta R^{AI,disp}$ captures the residual non-labor income not treated as a directly taxable capital/profit base. It may still finance purchases in formal markets and therefore induce consumption-tax revenue. Equivalently, one may define $\Delta M^{AI}=m^M(mpc^W\Delta W^{AI,disp}+mpc^\Pi\Delta\Pi^{AI,disp}+mpc^R\Delta R^{AI,disp})$ and subtract $\Delta M^{AI}$ once. The coefficient $m^M$ and an already-defined import increment must not both be applied to the same base. Consumption exemptions are assigned either to the taxable-consumption base through $(1-exempt^C)$ or to the construction of $\tau^{C,eff}$ through $\chi^C$, but not both. Robustness checks vary $mpc^W$, $mpc^\Pi$, $mpc^R$, import leakages, and VAT exemptions.

\subsection{Capital, foreign capture, and profit shifting}

AI-induced profit income may be domestic, foreign-owned, shifted, or deductible through payments for cloud services, licenses, software, or intellectual property. The object entering this block is the gross AI-induced profit base already separated from the national-accounts residual:
\begin{equation}
    \Delta \Pi^{AI,dom}=\chi^{dom}(1-\psi^{shift})\Delta \Pi^{AI,gross}.
\end{equation}
When the tax base rather than disposable income is required, the implementation uses $\Delta \Pi^{AI,dom}$ or a country-specific taxable-profit adjustment explicitly labelled as such.
Where withholding taxes or digital services taxes apply, foreign-remitted income can be added as a separate channel. This channel is treated as a robustness extension unless country-specific rules are observed.

\subsection{AI-specific rents and fiscal regimes}

AI-specific rents are written as
\begin{equation}
    Rent^{AI}=\lambda^{AI}g^{AI}Y^0.
\end{equation}
If $\tau^{AI}$ is statutory, then
\begin{equation}\label{eq:ai_eff_decomp_v21}
    \tau^{AI,eff}_{i,s,r,t}=\tau^{AI}_{i,r,t}\chi^{AI}_{i,s,r,t},
\end{equation}
where the AI-rent capturability term is explicitly decomposed as
\begin{equation}\label{eq:chi_ai_decomp_v21}
    \chi^{AI}_{i,s,r,t}
    =\chi^{AI,dom}_{i,s,t}(1-\psi^{shift,tech}_{i,s,t})
    \chi^{AI,compliance}_{i,r,t}\chi^{AI,enforce}_{i,r,t}.
\end{equation}
The AI-rent channel is interpreted as incremental rent capture not already included in the capital/profit taxable-base channel. Each simulation must declare one convention. If AI rents are inside the capital/profit base, Equation \eqref{eq:omega_k_net_v16} uses a unit-compatible $\lambda^{AI,Kbase}$ and $\omega^{K,net}=\max\{0,\omega^K-\lambda^{AI,Kbase}\}$ or sets $\lambda^{AI}\tau^{AI,eff}=0$. Let $\Delta\Pi^{AI,rent,gross}$ denote the AI-rent component of $\Delta\Pi^{AI,gross}$. Unit compatibility with the domestic post-shifting base used by $\omega^K$ requires
\begin{equation}
\begin{split}
    \lambda^{AI,Kbase}_{i,s,t}
    &=\chi^{dom}_{i,s,t}(1-\psi^{shift}_{i,s,t})
    \frac{\Delta\Pi^{AI,rent,gross}_{i,s,t}}{\Delta\Pi^{AI,gross}_{i,s,t}}
    \frac{\Delta\Pi^{AI,gross}_{i,s,t}/Y^0_{i,t}}{g^{AI}_{i,s,t}}\\
    &=\chi^{dom}_{i,s,t}(1-\psi^{shift}_{i,s,t})
    \frac{\Delta\Pi^{AI,rent,gross}_{i,s,t}/Y^0_{i,t}}{g^{AI}_{i,s,t}}.
\end{split}
\end{equation}
If $\lambda^{AI}$ is measured on the total AI shock, it must be rescaled to this domestic post-shifting base before it is subtracted from $\omega^K$. If AI rents are outside the measured capital/profit base, $\omega^{K,net}=\omega^K$ and the AI-rent channel is additional. This rule prevents double counting between the $K$ channel and the AI-specific rent channel.
Here $\chi^{AI,dom}$ captures the domestically reachable share of AI rents, $\psi^{shift,tech}$ captures technology-specific profit shifting through intellectual property, cloud location, licensing, royalties, or platform routing, and $\chi^{AI,compliance}$ and $\chi^{AI,enforce}$ capture legal compliance and administrative enforcement. The technology-shifting term is not assumed to equal ordinary capital-profit shifting; for emerging-market applications it may be larger. Existing digital services taxes in the region are typically turnover taxes on gross revenues rather than profit taxes; they map into this block as a higher effective rate on the rent margin with its own distortion costs. If $\tau^{AI,eff}$ is used directly, none of the components in $\chi^{AI}$ is multiplied again. Positive AI-rent taxation is not assumed in the status-quo baseline unless it exists in the country-year being modeled. The AI-rent channel must represent incremental capture outside the capital/profit channel. If AI rents are modeled as part of taxable profits, set $\lambda^{AI}\tau^{AI,eff}=0$ for the separate AI-rent term or net the corresponding component out of $\omega^K$. Technology profit shifting must be placed either inside $\chi^{AI}$ or inside residual leakage, but not both.

\subsection{Fixed costs, transition annualization, and marginal-equivalent deductions}

The implementation separates administrative flow costs, annualized one-off transition costs, fixed leakage, and marginal-equivalent deductions. All objects in the fiscal-space equation are annual shares of no-AI counterfactual GDP. Administrative costs are net of savings from programs that the new policy replaces:
\begin{equation}\label{eq:admin_net_costs_v21}
\begin{split}
    ac^{new,fix}_{i,p,s,r,t}&=\frac{AC^{new,fix}_{i,p,s,r,t}}{Y^0_{i,t}},\qquad
    ac^{exist,sav}_{i,p,t}=\frac{AC^{exist,sav}_{i,p,t}}{Y^0_{i,t}},\\
    ac^{net,fix}_{i,p,s,r,t}&=ac^{new,fix}_{i,p,s,r,t}-ac^{exist,sav}_{i,p,t}.
\end{split}
\end{equation}
The administrative-savings term is included only when the gross or net policy-cost object does not already include administrative expenditure from the replaced program. If $C^{exist}$ already includes both transfer spending and administrative spending, $ac^{exist,sav}$ is set to zero to avoid double counting. A scenario-pairing rule also applies: administrative savings presuppose that the replaced programs are actually terminated, which is the same replacement interpretation that defines $C^{exist}$ on the cost side. Pure-layering implementations, in which the new instrument is added on top of existing programs, must set $ac^{exist,sav}=0$. If $V^{gross}$ is reported with $ac^{exist,sav}>0$, the combination must be labelled a conservative hybrid: the numerator credits replacement administrative savings while the gross-cost denominator does not credit the replaced transfer budget.
Per-beneficiary administrative costs are systematically ordered: means-tested instruments exceed categorical instruments, which exceed universal instruments. Pilot implementations must not apply a uniform administrative-cost ratio across instruments, since that would compound the ideal-targeting bias in the GMI comparison.

Transition expenditure is compared with recurring annual policy costs only after annualization. Let $tr^{oneoff}_{i,s,r,t}=TR^{oneoff}_{i,s,r,t}/Y^0_{i,t}$ denote an implementation stock cost and $tr^{rec}_{i,s,r,t}$ any recurring transition-management cost. For horizon $H^{tr}$ and discount or financing rate $r^{tr}$, define
\begin{equation}\label{eq:transition_annualization_v21}
    a(H^{tr},r^{tr})=\begin{cases}
    \dfrac{r^{tr}(1+r^{tr})^{H^{tr}}}{(1+r^{tr})^{H^{tr}}-1}, & r^{tr}>0,\\[6pt]
    \dfrac{1}{H^{tr}}, & r^{tr}=0,
    \end{cases}
\end{equation}
\begin{equation}\label{eq:annualized_transition_v21}
    tr^{ann}_{i,s,r,t}=tr^{rec}_{i,s,r,t}+a(H^{tr}_{i,s,r},r^{tr}_{i,s,r})tr^{oneoff}_{i,s,r,t},
    \qquad
    leak^{fix}_{i,s,r,t}=\frac{Leak^{fix}_{i,s,r,t}}{Y^0_{i,t}}.
\end{equation}
A transition component already reported as an annual flow enters $tr^{rec}$ and is not annualized again. A one-off component enters only through $tr^{oneoff}$ and the annuity factor.

Marginal-equivalent deductions are per unit of the AI-induced fiscal-base shock:
\begin{equation}\label{eq:marginal_equiv_costs_v17}
    ac^{ME}_{i,s,r,t}\geq0,
    \qquad tr^{ME}_{i,s,r,t}\geq0,
    \qquad leak^{ME}_{i,s,r,t}\geq0.
\end{equation}
They are not computed by dividing fixed or annualized costs by $g^{AI}$. A cost component must be assigned to one margin only. If a cost is an annual administrative flow, an annualized one-off transition cost, or a fixed leakage, it enters $ac^{net,fix}$, $tr^{ann}$, or $leak^{fix}$. If it scales with the AI-induced output increment, it enters $ac^{ME}$, $tr^{ME}$, or $leak^{ME}$.

\subsection{Timing, base erosion, and behavioral response}

Lagged collection is modeled on the revenue object net of marginal-equivalent deductions but before fixed costs:
\begin{equation}
    rev^{gross}_{i,s,r,t}=\widetilde{MFC}^{gross}_{i,s,r,t}g^{AI,level}_{i,s,t}.
\end{equation}
The superscript gross here means before fixed costs only: $rev^{gross}$ is already net of marginal-equivalent deductions and therefore differs from $fs^{gross}=MFC^{gross}g^{AI}$, which is gross of them. Marginal-equivalent deductions are lagged together with the revenue they net out; fixed costs are not lagged.
\begin{equation}
    rev^{lag,gross}_{i,s,r,t}=\sum_{\ell=0}^{L}w^{lag}_\ell rev^{gross}_{i,s,r,t-\ell},
    \qquad \sum_{\ell=0}^{L}w^{lag}_\ell=1,
\end{equation}
where $rev^{gross}$ is already normalized by $Y^0$. Net lag-adjusted fiscal space is then
\begin{equation}
    fs^{eff,lag}_{i,p,s,r,t}=rev^{lag,gross}_{i,s,r,t}-ac^{net,fix}_{i,p,s,r,t}-tr^{ann}_{i,s,r,t}-leak^{fix}_{i,s,r,t}.
\end{equation}
Behavioral taxable-base erosion can be modeled as
\begin{equation}
    B^{j,taxable}=B^j(1-\epsilon^{ero,j}\tau^j),
\end{equation}
where $\epsilon^{ero,j}$ captures avoidance, relocation, deductions, or reporting responses not already embedded in the baseline capital-base chain. The erosion elasticities $\epsilon^{ero,j}$ are calibrated structural parameters in the sense of Section \ref{sec:parameters}: they carry no preferred point value, are varied over a declared grid or bounded support recorded in the calibration audit trail, and must satisfy $\epsilon^{ero,j}\tau^j<1$ so that taxable bases remain positive.
This extension is the mandatory robustness companion for reform-regime results ($r\geq1$).

\subsection{Fiscal-conversion parameter grid}

\begin{longtable}{>{\raggedright\arraybackslash}p{0.20\textwidth}>{\raggedright\arraybackslash}p{0.35\textwidth}>{\raggedright\arraybackslash}p{0.31\textwidth}}
\caption{Fiscal conversion robustness dimensions}\label{tab:fiscal_conversion_robustness}\\
\toprule
Parameter & Meaning & Robustness \\
\midrule
\endfirsthead
\toprule
Parameter & Meaning & Robustness \\
\midrule
\endhead
$\tau^{L,fisc,eff}$ & Effective labor-income capture & Country-specific; low/medium/high \\
$\tau^{K,fisc,eff}$ & Effective profit capture & Baseline: no extra ordinary profit shifting after $\Delta\Pi^{AI,dom}$; rate-side shifting only if $\psi^{shift}=0$ in the base chain \\
$\tau^{C,eff}$ & Effective consumption-tax capture & With and without VAT exemptions; do not also apply $(1-exempt^C)$ in the consumption base if exemptions enter here \\
$\tau^{AI,eff}$ & Effective AI-specific rent capture & Zero in status quo; positive scenarios \\
$\omega^L$ & Tax-base exposure weight for labor income & Derived from labor-share block \\
$\omega^K$ & Tax-base exposure weight for profits/capital & Labor-share block plus domestic-ownership/profit-shifting adjustment; $\Delta\Pi^{AI,dom}/Y^0$ divided by $g^{AI}$ \\
$\omega^C$ & Induced taxable-consumption response & Low/medium/high MPC, residual-income consumption, and import leakage \\
$mpc^R,\theta^{R,dom},\tau^{R,disp}$ & Residual non-labor income routed into consumption & Conservative/central/high residual-consumption pass-through \\
$\lambda^{AI}$ & AI-rent share of the shock, incremental to the capital/profit channel & Zero/low/medium/high; set to zero if included in $\omega^K$ \\
$MFC^{gross}$ & Drawn/calibrated gross marginal fiscal capture & Channel model, positive-capture conditional PERT, or unconditional two-component mixture \\
$MFC^{eff}$ & Derived effective reporting object after fixed costs and marginal-equivalent deductions & Computed from $fs^{eff}/g^{AI,level}$; not independently drawn \\
$ac^{net,fix},tr^{ann},leak^{fix}$; $ac^{ME},tr^{ME},leak^{ME}$ & Fixed GDP-share costs and marginal-equivalent deductions; each component assigned to one margin only & Low/medium/high \\
$\Omega^{res}$ & Residual stress-test penalty not already included elsewhere & Baseline zero; stress-test only \\
$AC^{new,fix},TR^{oneoff},Leak^{fix}$ & Nominal fixed-cost and transition-stock objects & Used only before GDP normalization to $Y^0$ \\
\bottomrule
\end{longtable}

The central reporting object remains $V^{gross}=fs^{eff}/c^{gross}$, with $V^{net}=fs^{eff}/c^{net}$ reported only as a robustness object when $c^{net}>0$; if $c^{net}=0$, the cell is labelled as already covered by existing spending.

\subsection{Historical fiscal-capture implementation details}

The historical plausibility check is implemented as a separate data object in the replication package. For each country-year, the replication file stores nominal GDP, tax revenue, total revenue, non-resource revenue where available, and government level. The preferred window is 2000--latest available year, with robustness windows 2010--latest and 2015--latest. Cleaning rules follow Section \ref{sec:historical_plausibility}. The historical output reports, for each country,
\begin{equation}
    \{P10,P25,P50,P75,P90\}\left(MFC^{hist,gross,+}_i\right),
    \qquad
    \{P10,P25,P50,P75,P90\}\left(Buoyancy^{hist}_i\right),
\end{equation}
and the share of years with negative gross capture.


\section{Informality and Fiscal-Capture Frictions}\label{app:informality_module}

This appendix formalizes how informality enters the fiscal-conversion block without becoming a generic penalty applied to the entire model. The key principle is that informality affects the capturability of specific tax bases. Labor informality, firm informality, consumption informality, exemptions, and profit shifting are therefore modeled separately. As in Appendix \ref{app:fiscal_conversion_module}, $\tau^{j,eff}$ abbreviates the fiscal-capture rates $\tau^{j,fisc,eff}$ of the main text.

\subsection{Informality by tax base}

The implementation distinguishes
\begin{equation}
    I^L_{i,t},\qquad I^F_{i,t},\qquad I^C_{i,t},
\end{equation}
where $I^L$ denotes labor informality, $I^F$ firm or business informality, and $I^C$ consumption or sales informality. In the baseline convention, profit shifting is represented by $\psi^{shift}_{i,s,t}$ only in the domestic profit-base chain, while consumption exemptions are represented by $exempt^C_{i,t}$.

\subsection{Taxable-share representation}

If a statutory or benchmark tax rate is used, capturability is used to construct the effective rate:
\begin{equation}\label{eq:tau_eff_informal_v10}
    \tau^{j,eff}_{i,s,r,t}=\tau^{j,base}_{i,r,t}\chi^j_{i,s,r,t},
    \qquad j\in\{L,K,C\}.
\end{equation}
The channel formula then uses $\tau^{j,eff}$ directly and does not multiply by $\chi^j$ again. Capturable shares may be constructed as
\begin{align}
    \chi^L_{i,s,r,t}&=(1-I^{L,marg}_{i,s,t})^{\rho_L},\\
    \chi^K_{i,s,r,t}&=(1-I^{F,marg}_{i,s,t})^{\rho_K},\\
    \chi^C_{i,s,r,t}&=(1-I^{C,marg}_{i,s,t})^{\rho_C}.
\end{align}
This is the baseline when profit shifting already enters $\Delta\Pi^{AI,dom}$ and $exempt^C$ already enters $\Delta C^{AI,taxable}$. If profit shifting is instead embedded in the effective capital-tax rate, use the alternative $\chi^K=(1-I^{F,marg})^{\rho_K}(1-\psi^{K,rate})^{\rho_\psi}$ and set $\psi^{shift}=0$ in the base-side capital chain. If consumption exemptions are instead embedded in the effective consumption-tax rate, use $\chi^C=(1-I^{C,marg})^{\rho_C}(1-exempt^C)^{\rho_E}$ and omit the $(1-exempt^C)$ multiplier from $\Delta C^{AI,taxable}$. If observed effective fiscal rates are used, additional informality, profit-shifting, or exemption penalties are reported only as robustness checks.

\subsection{Marginal informality of the AI shock}\label{subsec:marginal_informality}

If sectoral information is available, the marginal informality object is
\begin{equation}
\begin{aligned}
    I^{j,marg}_{i,s,t}&=\sum_k \omega^{AI}_{i,k,s,t}I^j_{i,k,t},
    &j&\in\{L,F,C\},\\
    \omega^{AI}_{i,k,s,t}&\geq0,
    &\sum_k\omega^{AI}_{i,k,s,t}&=1.
\end{aligned}
\end{equation}
The weights $\omega^{AI}_{i,k,s,t}$ are sectoral shares of the AI shock, so $I^{j,marg}_{i,s,t}$ is a weighted average in $[0,1]$. Where data allow, the weights are channel-matched: sectoral shares of the AI-induced labor-income increment for $j=L$, of the profit increment for $j=F$, and of induced consumption for $j=C$, with output-increment shares as the declared fallback. This prevents the model from mechanically applying aggregate national informality to shocks concentrated in more formal sectors such as finance, telecommunications, professional services, software, or formal manufacturing.

\subsection{No-double-counting rule}

\begin{longtable}{p{0.30\textwidth}p{0.23\textwidth}p{0.33\textwidth}}
\caption{No-double-counting rule for informality adjustments}\label{tab:no_double_counting_informality}\\
\toprule
Fiscal-rate source & Apply extra informality adjustment? & Reason \\
\midrule
\endfirsthead
\toprule
Fiscal-rate source & Apply extra informality adjustment? & Reason \\
\midrule
\endhead
Statutory or normative rate & Yes, inside $\tau^{eff}=\tau^{base}\chi$ & It does not embed compliance or informality \\
Observed effective rate from revenue/base & Not automatically & It already embeds average compliance, evasion, exemptions, and informality \\
Effective marginal rate for formal sector & Yes, with $I^{marg}$ if needed & The AI shock may occur partly outside the formal base \\
Aggregate tax-to-GDP ratio & No in the baseline & It already reflects the existing fiscal structure and informality \\
Labor channel, rate-side baseline & No base-side labor-informality penalty & Construct $\omega^L$ from gross $\Delta W^{AI}$ and include $\chi^L=(1-I^{L,marg})^{\rho_L}$ once in $\tau^{L,eff}$ \\
Labor channel, base-side alternative & No rate-side labor-informality penalty & Construct $\omega^L$ from $\Delta W^{taxable}$ and omit $\chi^L$ from $\tau^{L,eff}$ \\
Residual penalty $\Omega^{res}$ & Stress test only & Used only for frictions not included in channel-specific rates \\
\bottomrule
\end{longtable}
The baseline implementation derives effective fiscal rates from observed revenue composition, so no additional informality penalty is applied to them, consistent with the observed-effective-rate row of Table \ref{tab:no_double_counting_informality}.
Legal consumption exemptions and consumption informality overlap empirically: informal purchases concentrate in exempt staples. To avoid discounting the same untaxed consumption twice, $exempt^C$ must be measured on the formal consumption basket, or $I^C$ on the taxable basket, and the calibration registry states which convention is used. A parallel overlap applies to the capital channel: mixed income and informal-business income already excluded from the taxable base through $\chi^{Kbase}$ must not be penalized again through $I^{F,marg}$. Under the baseline, $\chi^{Kbase}$ is a composition filter selecting the profit-type components of the non-labor residual, while $I^{F,marg}$ is a compliance filter on firms inside that taxable-profit perimeter; the calibration registry states that $I^{F,marg}$ is measured on the perimeter-consistent firm population, or assigns the informal-business margin to one side only.

\subsection{Informality in adoption versus fiscal capture}

Informality can constrain both adoption and fiscal capture, but through different equations. In adoption, informality affects $q^{prod}$ and therefore $g^{AI}$. In fiscal conversion, informality affects $\tau^{j,eff}$ and therefore $fs^{eff}$ conditional on $g^{AI}$. The calibration registry declares whether an informality variable is being used on the adoption side, the fiscal side, or both; if both, the interpretation must be that they represent separate mechanisms rather than a repeated penalty on the same object.

\subsection{Formalization-enhancing AI as a scenario}

The baseline does not assume that AI automatically formalizes the economy. A reform scenario may allow digitalization, e-invoicing, payments traceability, or AI-assisted tax administration to reduce informality:
\begin{equation}
    I^{j,post}_{i,s,r,t}=I^j_{i,t}\left(1-\varphi^j_{r,t}q^{prod}_{i,s,t}\right).
\end{equation}
For each labor, firm, and consumption informality channel, the object is truncated to keep $I^{j,post}\in[0,1]$. This extension is inactive in the status quo and activated only in improved-compliance or combined-reform scenarios.
The endogenous formalization term applies only to fiscal capturability within the period: $I^{j,post}$ enters $\chi^j$ and $\tau^{j,eff}$, while the adoption-side informality in $\ell^q$ and $\bar q$ is held at observed values to avoid a simultaneous loop between adoption and the adoption ceiling. A dynamic extension may feed $I^{j,post}$ back into adoption with a one-period lag; it must then be labelled as a dynamic-feedback robustness scenario.

The symmetric adverse scenario is displacement-driven informalization: when automation pressure is strong and reinstatement weak, displaced formal workers may move into informality. A labelled robustness scenario allows $I^{L,post}_{i,s,r,t}=I^L_{i,t}\left(1+\varphi^{disp}_{r,t}(\delta_s-\zeta_s RST_{i,t})E^{disp}_{i,t}q^{prod}_{i,s,t}\right)$, truncated to keep $I^{L,post}\in[0,1]$, with the same fiscal-side-only timing convention as the formalization scenario. The baseline holds informality fixed; of the two endogenous deviations, only reporting the favorable one would bias automation-heavy scenarios optimistically, so implementations that activate formalization scenarios must also report this adverse counterpart.

\subsection{Elasticity and robustness grid}

The baseline elasticities $\rho_L,\rho_K,\rho_C,\rho_E,\rho_\psi$ are varied over the grid $\{0.5,1,1.5,2\}$. The profit-shifting elasticity $\rho_\psi$ is active only in the alternative rate-side convention for capital taxation. The exemption elasticity $\rho_E$ is active only in the alternative convention where consumption exemptions are embedded in $\chi^C$ rather than in $\Delta C^{AI,taxable}$. They are not treated as known structural constants.


\section{Threshold Inversion Module}\label{app:threshold_module}

This appendix formalizes additional threshold inversions used as stress tests. All thresholds are algebraic inversions of the same feasibility condition, not independent assumptions and not forecasts.

\subsection{Notation and no-double-counting rule}

Let
\begin{equation}
    R^{req}_{i,p,t}(\xi)=(1+\xi)c^{gross}_{i,p,t},
    \qquad
    R^{req,DC}_{i,p,t}(\xi)=(1+\xi)c^{gross}_{i,p,t}+s^{PB}_{i,t}.
\end{equation}
If the threshold uses gross fiscal conversion, the required gross revenue includes only fixed-cost subtractions. Marginal-equivalent deductions enter the denominator through $\widetilde{MFC}^{gross}$ and must not be added again to the required-revenue term:
\begin{equation}
    R^{req,gross}_{i,p,s,r,t}(\xi)=R^{req}_{i,p,t}(\xi)+ac^{net,fix}_{i,p,s,r,t}+tr^{ann}_{i,s,r,t}+leak^{fix}_{i,s,r,t},
\end{equation}
\begin{equation}
    R^{req,gross,DC}_{i,p,s,r,t}(\xi)=R^{req,DC}_{i,p,t}(\xi)+ac^{net,fix}_{i,p,s,r,t}+tr^{ann}_{i,s,r,t}+leak^{fix}_{i,s,r,t}.
\end{equation}
The implementation must not add $ac^{net,fix}$, $tr^{ann}$, or $leak^{fix}$ to a threshold that already uses $MFC^{eff}$. It must also not add $ac^{ME}$, $tr^{ME}$, or $leak^{ME}$ to $R^{req,gross}$, because those terms are already embedded in $\widetilde{MFC}^{gross}=MFC^{gross}-ac^{ME}-tr^{ME}-leak^{ME}$. If $\Omega^{res}$ is activated, it enters only as a stress-test penalty: either use $fs^{eff,res}=fs^{eff}-\Omega^{res}$ or add $\Omega^{res}$ to the required-revenue object, but not both.

The closed-form threshold formulas in this appendix are valid for the historical reduced-form mode or for channel-based calculations in which $\omega^j$, $\lambda^{AI}$, and $\widetilde{MFC}^{gross}$ are held fixed at the evaluated scenario. If the channel bases $\Delta B^{j,tax}$ are recomputed as $g^{AI}$, $T$, or $q^{prod}$ changes, the implementation solves the threshold directly from $fs^{eff,ch}(\cdot)$ in Equation \eqref{eq:channel_revenue_function_v24}.

\subsection{Minimum AI-rent capture threshold}

Suppose the gross fiscal conversion block is written with the AI-rent term separated from the labor, capital, and consumption channels, $MFC^{base,gross}_{i,s,r,t}=MFC^{gross}_{i,s,r,t}-\lambda^{AI}_{i,s,t}\tau^{AI,eff}_{i,s,r,t}$ in the notation of Equation \eqref{eq:mfc_gross_channel_v10}, so that
\begin{equation}
    fs^{eff}_{i,p,s,r,t}=g^{AI}_{i,s,t}\left[\widetilde{MFC}^{base,gross}_{i,s,r,t}+\tau^{AI,eff}_{i,s,r,t}\lambda^{AI}_{i,s,t}\right]-ac^{net,fix}_{i,p,s,r,t}-tr^{ann}_{i,s,r,t}-leak^{fix}_{i,s,r,t},
\end{equation}
where
\begin{equation}
    \widetilde{MFC}^{base,gross}_{i,s,r,t}=MFC^{base,gross}_{i,s,r,t}-ac^{ME}_{i,s,r,t}-tr^{ME}_{i,s,r,t}-leak^{ME}_{i,s,r,t}.
\end{equation}
Solving for the minimum effective AI-rent tax gives
\begin{equation}\label{eq:tau_ai_eff_star_v10}
    \tau^{AI,eff,*}_{i,p,s,r,t}(\xi)=
    \max\left\{0,
    \frac{\frac{R^{req,gross}_{i,p,s,r,t}(\xi)}{g^{AI}_{i,s,t}}-\widetilde{MFC}^{base,gross}_{i,s,r,t}}
    {\lambda^{AI}_{i,s,t}}
    \right\},
\end{equation}
provided $g^{AI}_{i,s,t}>0$ and $\lambda^{AI}_{i,s,t}>0$. If the policy variable is a statutory rate and $\tau^{AI,eff}=\tau^{AI}\chi^{AI}$, then for $\chi^{AI}_{i,s,r,t}>0$,
\begin{equation}\label{eq:tau_ai_stat_star_v10}
    \tau^{AI,*}_{i,p,s,r,t}(\xi)=\frac{\tau^{AI,eff,*}_{i,p,s,r,t}(\xi)}{\chi^{AI}_{i,s,r,t}}.
\end{equation}
If $\chi^{AI}_{i,s,r,t}=0$ and $\tau^{AI,eff,*}>0$, the cell is AI-rent-capture-impossible. If $\tau^{AI,eff,*}>\bar\tau^{AI,eff}$ or $\tau^{AI,*}>\bar\tau^{AI}$, the policy-scenario pair is infeasible through the AI-rent-tax channel.

\subsection{Minimum productive-adoption threshold}

To avoid circularity and relative-sample inflation, the adoption threshold uses a fixed external reference score $S^{ref}_{s,t}=S^{frontier}_{s,t}$ that does not depend on the country-policy cell being inverted or on the internal pilot sample. Under the baseline no-double-counting score,
\begin{equation}
    S^{NDC}_{i,s,t}=\mathbf{1}\{q^{prod}_{i,s,t}>0,\ E^{prod}_{i,t}>0\}(\tilde E^{prod}_{i,t})^\beta(\tilde q^{prod,num}_{i,s,t})^\gamma,
\end{equation}
The closed-form translation and adoption thresholds, including the conversion back from the floor-regularized scale and lower-support truncation, are stated in Proposition \ref{prop:q_threshold_v10}. If $\phi^{Y,nom}_s\leq0$ or $\widetilde{MFC}^{gross}\leq0$ with positive required revenue, no finite adoption threshold exists and the cell is labelled non-positive technology/fiscal support. If $T^{req,static}>1$, even the saturated value $T=1$ cannot cross the threshold and the cell is classified as technology/fiscal-support insufficient. Diagnostics apply to the truncated original-scale object $q^{prod,*}$; if $q^{prod,*}>\bar q_i$, feasibility cannot be obtained by adoption alone under the maintained exposure, AI-productivity, and fiscal-conversion assumptions. The unique operational definition of floor-driven is the one in Section \ref{subsec:impossibility_score}, Equation \eqref{eq:floor_driven_operational_v25}: $q^{prod,*}=0$, $\epsilon>0$, and the crossing fails when $\epsilon=0$. Such a cell is not counted as adoption-feasible. A substantive adoption crossing must satisfy $q^{prod,*}>0$ and $q^{prod,*}\leq\bar q_i$.

\subsection{Minimum base-shock threshold}

Under the one-period static approximation and the frozen-weight convention, the required nominal GDP-equivalent base shock is
\begin{equation}\label{eq:phi_star_static_app_v25}
    \phi^{*,static}_{i,p,s,r,t}(\xi)
    =\frac{R^{req,gross}_{i,p,s,r,t}(\xi)}
    {T_{i,s,t}\widetilde{MFC}^{gross}_{i,s,r,t}}.
\end{equation}
This threshold is valid when $T_{i,s,t}>0$ and $\widetilde{MFC}^{gross}_{i,s,r,t}>0$, under the same existence and frozen-weight rules as $T^*$. For accumulated horizons, $\phi^*$ is solved numerically from Equation \eqref{eq:accumulated_threshold_condition_v22}. If $\phi^{*,static}_{i,p,s,r,t}(\xi)>\phi^{Y,nom,max}_s$, the cell is technology-scenario-insufficient.

\subsection{Existence conditions for adoption thresholds}\label{subsec:adoption_threshold_existence}

An interior adoption threshold exists only if $R>0$, $\phi^{Y,nom}_s>0$, $\widetilde{MFC}^{gross}>0$, $\gamma>0$, $E^{prod}>0$, $\tilde E^{prod}>0$, and $S^{ref}=S^{frontier}>0$. This list is identical to the condition $Valid^{T*}_{i,p,s,r,t}(\xi)=1$ in Proposition \ref{prop:q_threshold_v10}. A growth threshold similarly requires $\widetilde{MFC}^{gross}>0$, and a gross-capture threshold requires $g^{AI}>0$. If $\widetilde{MFC}^{gross}\leq0$ while required revenue is positive, the draw is kept as an economic failure draw but no finite positive growth or adoption threshold is reported. In simulations, any use of $\epsilon>0$ is numerical regularization conditional on positive adoption; substantive feasibility at zero adoption or zero productive exposure is ruled out by construction through the gate $q^{prod}\leq0\ \text{or}\ E^{prod}\leq0\Rightarrow S=T=g^{AI}=0$.

\subsection{Preparedness and exposure thresholds}

Preparedness and exposure thresholds are reported as stress tests only. If $q^{prod}$ is simulated from preparedness, a preparedness threshold can be computed from the adoption equation. If preparedness is also in $S^{full}$, then the specification must be labeled full-score robustness rather than no-double-counting baseline.

\subsection{Impossibility diagnostics}

The threshold module reports realized coverage, support headroom, and absolute gaps separately. Realized coverage is
\begin{align}\label{eq:distance_ratios_main_v19}
    D^{g,real}&=\frac{g^{AI}}{g^{AI,*}},
    &D^{MFC,gross,real}&=\frac{MFC^{gross}}{MFC^{gross,*}},\\
    D^{q,real}&=\frac{q^{prod}}{q^{prod,*}}\quad\text{if }q^{prod,*}>0,
    &D^{AI,eff,real}&=\frac{\tau^{AI,eff}}{\tau^{AI,eff,*}}\quad\text{if }\tau^{AI,eff,*}>0.
\end{align}
Support headroom is
\begin{align}\label{eq:headroom_ratios_main_v19}
    H^g&=\frac{g^{AI,max}}{g^{AI,*}},
    &H^q&=\frac{\bar q_i}{q^{prod,*}}\quad\text{if }q^{prod,*}>0,\\
    H^{AI,eff}&=\frac{\bar\tau^{AI,eff}}{\tau^{AI,eff,*}}\quad\text{if }\tau^{AI,eff,*}>0.
\end{align}
Values above one mean that the realized value or support ceiling covers the requirement within its family. The effective AI-rent bound combines the statutory cap and capturability, $\bar\tau^{AI,eff}_{i,s,r,t}=\bar\tau^{AI}_{i,r,t}\chi^{AI}_{i,s,r,t}$. Realized and headroom families are never pooled into one cross-margin ranking. A derived effective-conversion distance must be separately labelled $D^{MFC,eff,real}=MFC^{eff}/MFC^{eff,*}$ and used only for derived reporting; gross realized capture remains the baseline plausibility metric because historical capture is gross.
If $q^{prod,*}=0$, neither $D^{q,real}$ nor $H^q$ is reported. Adoption-margin already satisfied requires feasibility without reliance on the numerical floor and without violating the zero-adoption/zero-exposure gate; otherwise the cell is floor-driven or non-substantive and not adoption-feasible. If $0<q^{prod,*}\leq\bar q_i$, support is sufficient; if $q^{prod,*}>\bar q_i$, the ceiling fails. The absolute excess $Gap^q=q^{prod,*}-\bar q_i$ is reported as a gap, never a distance ratio.


For accumulated-horizon implementations, the scenario-supported maximum AI level gain is
\begin{equation}
    g^{AI,max}_{s,H}=\prod_{\tau=1}^{H}\left(1+\phi^{Y,nom,max}_{s,\tau}D_{s,\tau}T^{max}\right)-1.
\end{equation}
The static object $g^{AI,max}_{s,t}=\phi^{Y,nom,max}_sD_{s,t}T^{max}$ is used only for one-period static runs. The consolidated classification of impossibility diagnostics is in Table \ref{tab:impossibility_diagnostics_v10}; this module reports the ratios and gaps defined above.

\subsection{Comparative statics}

Under the frozen-weight convention, each required threshold is increasing in policy cost and decreasing in the realized values of the other margins: fiscal capture for $g^{AI,*}$; exposure and adoption for thresholds that operate through $T$; and the productivity scenario for $MFC^{gross,*}$. A threshold does not depend on its own realized value. These comparative statics are algebraic and should be reported as threshold movements rather than causal estimates.

\subsection{Reporting protocol}

Every threshold table should state whether it uses gross or effective fiscal conversion, gross or net policy cost, debt-consistent or non-debt requirements, and whether the adoption score is no-double-counting or full-score robustness.

\subsection{Pressure-without-capacity diagnostic}

The main text refers to the following optional ranking diagnostic.

The optional pressure-without-capacity diagnostic is defined as
\begin{equation}\label{eq:pwc_index_v22}
    PWC_{i,p,s,r,t}=E^{disp}_{i,t}(1-A_{i,t})I^L_{i,t}\max\{0,1-V^{gross}_{i,p,s,r,t}\},
\end{equation}
where $E^{disp}$ captures displacement exposure, $(1-A)$ captures low preparedness, $I^L$ captures labour informality, and $\max\{0,1-V^{gross}\}$ captures uncovered fiscal distance. The index is a ranking diagnostic only; it is not a feasibility threshold and is not added to the fiscal-space equation.


\section{Uncertainty Propagation and Robustness Module}\label{app:uncertainty_module}

This appendix formalizes the robustness and Monte Carlo protocol. The Monte Carlo exercise propagates parameter uncertainty through the fiscal-feasibility framework; it is not a forecast of AI development.


\subsection{Monte Carlo inputs}

The uncertain-input vector is mode-specific. In the channel-based fiscal mode, draw channel rates and construct $MFC^{gross}$ from them. In the historical reduced-form mode, draw $MFC^{gross}$ directly from the historical gross-capture distribution and do not also use channel rates to construct it. The positive-capture PERT version is a conditional-feasibility benchmark; the two-component mixture is the unconditional fiscal-risk benchmark.

The implementation must also declare its adoption mode. In observed-adoption mode, $q^{prod}$ is a primitive drawn from observed adoption or measurement-error distributions. In the baseline anchored simulated-adoption mode, $A,I,Gap,\lambda_q,\nu_A,\nu_I,\nu_G$, and $q^{use,target}_{i,t_0}$ are primitives, while $\mu_{i,s,t_0}$ is derived by the target inversion. In structural-intercept robustness mode, $\mu_s$ or $\mu_s+\alpha_i$ is drawn directly and the target inversion is not used. In all simulated modes, $q^{use}$ and $q^{prod}$ are derived endogenously.

Channel-based fiscal mode:
\begin{equation}
\begin{split}
    \Theta^{(m)}_{channel}=\{&\phi^{raw}_{s,u},\kappa^{u\to Y},\pi^{AI,Y},A,E^{prod},I,Gap,\lambda_q,\nu_A,\nu_I,\nu_G,\omega_I,\omega_G,q^{use,target},\\
    &\tau^{L,fisc,eff},\tau^{K,fisc,eff},\tau^{C,eff},\lambda^{AI},\tau^{AI},\chi^{AI,dom},\psi^{shift,tech},\chi^{AI,compliance},\chi^{AI,enforce},\\
    &ac^{net,fix},tr^{ann},leak^{fix},ac^{ME},tr^{ME},leak^{ME},\\
    &\rho_L,\rho_K,\rho_C,\rho_E,\rho_\psi,\alpha,\beta,\gamma,
    LS,E^{auto},E^{aug},RST,\delta,\psi,\zeta,\lambda_{LS},\chi^{Kbase},\\
    &mpc^W,mpc^{\Pi},mpc^R,m^M,exempt^C,\chi^{dom},\psi^{shift},\theta^{R,dom},\tau^{L,disp},\tau^{K,disp},\tau^{R,disp},\Omega^{res},s^{PB},c^{gross},c^{net}\}^{(m)}.
\end{split}
\end{equation}
Historical reduced-form fiscal mode:
\begin{equation}
\begin{split}
    \Theta^{(m)}_{hist}=\{&\phi^{raw}_{s,u},\kappa^{u\to Y},\pi^{AI,Y},A,E^{prod},I,Gap,\lambda_q,\nu_A,\nu_I,\nu_G,q^{use,target},\omega_I,\omega_G,
    MFC^{gross},ac^{net,fix},tr^{ann},leak^{fix},\\
    &ac^{ME},tr^{ME},leak^{ME},
    \alpha,\beta,\gamma,\Omega^{res},s^{PB},c^{gross},c^{net}\}^{(m)}.
\end{split}
\end{equation}
Nominal $AC$, $TR$, and $Leak$ objects are reporting objects derived from $Y^0$ when needed. They are not drawn independently from normalized fixed-cost or marginal-equivalent objects in the same simulation. Observed variables may be fixed in the baseline, but the replication table must label them as fixed and give source-year values. If they are stochastic, the table must give distribution, support, source, and scenario cell.

\begin{footnotesize}
\begin{longtable}{>{\raggedright\arraybackslash}p{0.18\textwidth}>{\raggedright\arraybackslash}p{0.27\textwidth}>{\raggedright\arraybackslash}p{0.28\textwidth}>{\raggedright\arraybackslash}p{0.17\textwidth}}
\caption{Monte Carlo input protocol}\label{tab:mc_inputs}\\
\toprule
Parameter block & Meaning & Baseline distribution & Robustness \\
\midrule
\endfirsthead
\toprule
Parameter block & Meaning & Baseline distribution & Robustness \\
\midrule
\endhead
AI shock block & Raw AI shock, output bridge, nominal bridge, and accumulated level object & Bounded PERT on raw support; bridge fixed or PERT & Triangular or truncated-uniform shock supports \\
Adoption block & $A$, $E^{prod}$, $I$, $Gap$, $q^{prod}$, $\lambda_q$, $\omega_I$, $\omega_G$, and anchored target moment & Observed inputs fixed or beta/PERT measurement error; simulated adoption derives $q^{use}$ and $q^{prod}$ from the anchored diffusion module & Slow/fast diffusion and structural-intercept robustness \\
Fiscal-rate block & Labor, capital, consumption, and AI-rent fiscal rates & Observed effective rates where available; otherwise PERT or fixed scenario values & Low/high fiscal capture \\
AI-rent capture block & $\lambda^{AI}$, statutory $\tau^{AI}$, and the components of $\chi^{AI}$: domestic reach, technology profit shifting, compliance, and enforcement & Zero AI-rent tax in status quo; positive scenarios use bounded shares & No AI-rent tax; high technology-shifting stress \\
Cost and leakage block & Net administrative cost, annualized transition cost, fixed leakage, and marginal-equivalent deductions & Bounded PERT or fixed source values; each cost assigned to exactly one margin & High-cost/high-leakage stress \\
Tax-base elasticities & Informality elasticities; $\rho_E$ only under the labelled exemptions-in-rate convention; $\rho_\psi$ only under labelled rate-side profit shifting & Bounded PERT; optional terms active only in their labelled conventions & Discrete grid \\
Behavioral erosion & $\epsilon^{ero,j}$ taxable-base erosion under reform regimes & Fixed at zero in the status-quo baseline & Mandatory declared grid or bounded PERT companion for $r\geq1$ \\
Residual penalty & $\Omega^{res}$ for residual stress-test frictions not counted elsewhere & Fixed at zero in baseline & Positive stress-test draws only \\
Translation elasticities & $\alpha$, $\beta$, and $\gamma$ & Discrete grid & Factorial grid \\
Labor-share block & Automation, augmentation, reskilling, and labor-share damping & Bounded PERT & Truncated-uniform robustness \\
Policy-cost block & Gross cost, net cost, and debt-consistent guardrail & Fixed baseline or scenario objects & Net-cost and debt-guardrail robustness \\
\bottomrule
\end{longtable}
\end{footnotesize}

\subsection{Scenario-consistent dependence}

The baseline does not require all parameters to be independent. Correlated latent shocks can encode $Corr(A,q)>0$, $Corr(I,q)<0$, $Corr(I,MFC)<0$, $Corr(A,MFC)>0$, and additionally $Corr(I,\chi^{Kbase})<0$ (informality raises the mixed-income share of the non-labor residual) and $Corr(A,E^{prod})>0$ (more prepared economies are also more exposed). Because the signs cluster advantages, scenario-consistent dependence widens the cross-country dispersion of $V$ relative to independence: strong countries cross earlier and weak countries later. These are robustness assumptions, not estimated causal correlations.

\subsection{Monte Carlo outputs}

For each country-policy-scenario-regime cell, report
\begin{equation}
    \{P5,P25,P50,P75,P95\}(V^{gross})
\end{equation}
and, when relevant,
\begin{equation}
    \{P5,P25,P50,P75,P95\}(V^{net})\quad\text{only for cells with }c^{net}>0.
\end{equation}
The crossing-probability grid is
\begin{equation}
    \widehat{\pi}(\xi)=\frac{1}{M}\sum_{m=1}^{M}\mathbf{1}\{V^{gross,(m)}\geq1+\xi\},
    \qquad \xi\in\{0,0.05,0.10,0.25,0.50\}.
\end{equation}
The basic-valid and support-valid versions are
\begin{equation}
    \widehat{\pi}^{basic}(\xi)=
    \frac{\sum_m Valid^{basic,(m)}\mathbf{1}\{V^{gross,(m)}\geq1+\xi\}}
    {\sum_m Valid^{basic,(m)}}
    \quad\text{if }\sum_m Valid^{basic,(m)}>0,
\end{equation}
\begin{equation}
    \widehat{\pi}^{support}(\xi)=
    \frac{\sum_m Valid^{support,(m)}\mathbf{1}\{V^{gross,(m)}\geq1+\xi\}}
    {\sum_m Valid^{support,(m)}}
    \quad\text{if }\sum_m Valid^{support,(m)}>0.
\end{equation}
Report $\pi^{threshold|basic}$ and $IS^{support|threshold}$ to separate non-computable thresholds from support violations, and retain $\pi^{support|basic}$ as the aggregate support-valid share. The debt-consistent version additionally requires $fs^{eff,(m)}\geq(1+\xi)c^{gross,(m)}+s^{PB,(m)}$.

\begin{longtable}{p{0.18\textwidth}p{0.36\textwidth}p{0.34\textwidth}}
\caption{Required Monte Carlo output layers}\label{tab:mc_output_layers}\\
\toprule
Layer & Output & Purpose \\
\midrule
\endfirsthead
\toprule
Layer & Output & Purpose \\
\midrule
\endhead
Distribution & $P5$, $P25$, $P50$, $P75$, $P95$ of $V^{gross}$ & Shows downside, median, and upside feasibility \\
Crossing probability & $\widehat{\pi}(\xi)$, $\widehat{\pi}^{basic}(\xi)$, and $\widehat{\pi}^{support}(\xi)$ & Shows whether feasibility survives buffers with all-draw, basic-valid, and support-valid denominators \\
Debt-consistent probability & $\widehat{\pi}^{DC}(\xi)$ & Separates accounting coverage from fiscal sustainability \\
Failure severity & $E[1+\xi-V\mid V<1+\xi]$ & Distinguishes slight failure from severe failure \\
Validity / support chain & $\pi^{basic}$, $\pi^{threshold|basic}$, and $IS^{support|threshold}$ & Separates denominator-invalid draws, non-computable thresholds, and bounded-support violations \\
Driver ranking & Spearman/tornado/Sobol outputs & Identifies what moves $V$ \\
\bottomrule
\end{longtable}

\subsection{Global sensitivity analysis}

A parsimonious first implementation reports Spearman rank correlations
\begin{equation}
    \rho^S_j=Corr_S(\theta_j,V^{gross}).
\end{equation}
A tornado chart varies one input at a time between low and high values. Where inputs are independent, first-order and total-order Sobol indices may be reported \citep{morris1991factorial,saltelli2008global,sobol2001global}:
\begin{equation}
    S_j=\frac{Var_{\theta_j}(E[V\mid\theta_j])}{Var(V)},
    \qquad
    S^T_j=1-\frac{Var_{\theta_{-j}}(E[V\mid\theta_{-j}])}{Var(V)}.
\end{equation}
If the Monte Carlo imposes dependence, Sobol indices are reported only in an independent diagnostic run or with a dependence-aware method.

\subsection{Deterministic stress tests}

The deterministic stress menu spans low AI productivity, slow adoption, weak fiscal capture, high informality, high administrative cost, high leakage, no AI-rent tax, no formalization effect, and no labor augmentation; the executed stress and ablation outputs are recorded in the replication package.

\subsection{Ablation, placebo, and negative-control outputs}

The replication package includes temporal placebo, sectoral negative-control, leave-one-source-out, and channel-ablation outputs.

\subsection{Numerical convergence and reproducibility}

The replication package reports the random seed, draw count, convergence checks for crossing probabilities (its Monte Carlo convergence output), support-valid draw share, code version, and data version. Any table using $V^{net}$ must also report $V^{gross}$.

\section{Calibration Matrix and Assumption Ledger}\label{app:calibration_ledger}

This appendix consolidates the minimum parameter-reporting standard and the ledger of maintained assumptions with their bias directions.

\subsection{Master parameter-assignment and calibration matrix}

The following matrix is the minimum reporting standard for the empirical application. It is designed to make each value reproducible and to prevent the model from hiding fragile assumptions inside a single composite index.

\scriptsize
\begin{longtable}{p{0.14\textwidth}p{0.22\textwidth}p{0.25\textwidth}p{0.20\textwidth}p{0.11\textwidth}}
\caption{Master parameter-assignment and calibration matrix}\label{tab:master_calibration_matrix}\\
\toprule
Input & Definition and unit & Baseline assignment & Uncertainty treatment & Primary source discipline \\
\midrule
\endfirsthead
\toprule
Input & Definition and unit & Baseline assignment & Uncertainty treatment & Primary source discipline \\
\midrule
\endhead
$Y_{i,t}$, $Pop_{i,t}$, $N_{i,p,t}$ & GDP, population, and eligible population for policy $p$ & Observed country-year values; keep nominal and PPP versions separate & Fixed in baseline; measurement-error check if data are interpolated & World Bank PIP/WDI; national accounts \\
$c^{gross,0}_{i,p,t}$, $c^{net,0}_{i,p,t}$ & Annual gross or net policy cost as share of no-AI counterfactual GDP $Y^0_{i,t}$ & Construct from benefit rule $B_{i,p,t}$ and eligible population; subtract existing spending only in the net-cost specification; post-AI denominators are robustness only & Low/central/high benefit rules; no mixing gross and net cost & ASPIRE, PIP, national budgets, \citet{gentilini2020ubi} \\
$\eta_{MUT},\eta_{PBI},\eta_{UBI},\chi_{GMI},B^{PEN},\theta^{target}$ & Benefit-rule parameters, pension benefit level, and GMI targeting-and-behavior loading & Indicative support $[0.25,1.00]$ for poverty-line fractions and gap coverage, including $B^{PEN}/PL$; $\theta^{target}\in[1.0,1.5]$ for labelled GMI robustness; finalized in the baseline calibration registry (replication package, appendix\_i files) & Low/central/high grid or bounded PERT within declared support & PIP, ASPIRE, and national benefit rules \\
$\phi^{raw}_{s,u}$, $\phi^{Y,nom}_s$, $g^{AI,level}$ & Raw source-unit AI shock, nominal GDP-equivalent annual bridge, and accumulated level shock used against recurring costs & Draw raw support from Table \ref{tab:phi_numerical_closure}; convert through $\kappa^{u\to Y}$ and $\pi^{AI,Y}$; accumulate before feasibility comparison & Bounded PERT on raw support; bridge uncertainty reported separately; triangular/uniform only robustness & \citet{acemoglu2024simple}, \citet{filippucci2024miracle}, \citet{oecd2025g7} \\
$A_{i,t}$ & AI preparedness, normalized $[0,1]$ & IMF AIPI or documented reconstruction from its four dimensions: digital infrastructure, human capital and labour-market policies, innovation/economic integration, regulation/ethics & Fixed observed index plus normalization robustness; avoid using as causal estimate & \citet{imf2024aipi,cazzaniga2024genai} \\
$Gap_{i,t}$ & Digital and skills gap, normalized to $[0,1]$ & Construct only from indicators excluded from AIPI & Fixed observed index plus normalization robustness & ITU and World Bank digital indicators \\
$E^{prod}_{i,t}$, $\varpi$ & Productive AI exposure and occupational-versus-sectoral combination weight & Prefer occupation/sector exposure mapped to employment or value added; for LAC, use productivity-enhancing exposure as baseline rather than total exposure & Beta/PERT distribution around low-central-high exposure; $\varpi\in\{0.25,0.50,0.75\}$ or a declared weight; broad exposure is robustness, not baseline & \citet{worldbankilo2024lacgenai,cazzaniga2024genai} \\
$E^{disp}_{i,t}$ & Displacement or automation exposure & Separate from $E^{prod}$; used in pressure and labour-share modules, not directly as productivity opportunity & Scenario range; never add to $E^{prod}$ & \citet{worldbankilo2024lacgenai,cazzaniga2024genai} \\
$q^{use}_{i,s,t}$, $q^{prod}_{i,s,t}$, $\omega_I$, $\omega_G$ & AI use, productive adoption, and adoption-ceiling parameters; adoption shares lie in $[0,1]$ & If observed, use firm/worker adoption data; if simulated, use logistic module with country ceiling $\bar q_i$ & Logit-normal or bounded PERT for adoption; grid or bounded PERT with declared support for $\omega_I$ and $\omega_G$; enforce $0\le q^{prod}\le\bar q_i\le1$; impose $q^{prod}\le q^{use}$ only under non-decreasing-use robustness & Diffusion literature \citep{bass1969newproduct,rogers2003diffusion}; micro-productivity evidence \citep{noy2023experimental,brynjolfsson2023generative} \\
$\lambda_q$ & Productive-adoption lag parameter & Baseline grid $\{0.15,0.30,0.50\}$; $1.00$ only as instant-adoption stress & Discrete grid & Diffusion theory and adoption-lag logic \\
$\nu_A,\nu_I,\nu_G,q^{use,target}$ & Logistic adoption coefficients and anchored adoption moment & Calibrate the anchored simulated-adoption mode in Appendix \ref{app:adoption_module} & Grid or bounded PERT with declared support & Appendix \ref{app:adoption_module}; diffusion literature; national ICT/firm surveys or World Bank technology adoption surveys for the target moment \\
$S^{frontier}_{s,t}$ & External translation benchmark & Document the frontier-economy set and benchmark year & Fixed within each threshold inversion; alternative external sets as robustness & External frontier sample; Appendix \ref{app:ai_shock_module} \\
$\alpha,\beta,\gamma$ & Elasticities in the translation factor & $\beta=\gamma=1$ baseline; $\alpha=0$ under the no-double-counting baseline ($\alpha=1$ only in full-score robustness; see Appendix \ref{app:ai_shock_module} elasticity grid) & Grid $\{0.5,1,1.5,2\}$ for active elasticities and global sensitivity & Model discipline; no direct causal claim \\
$RST_{i,t}$, $\delta_s,\psi_s,\zeta_s,\lambda_{LS}$ & Reskilling/task-reinstatement capacity in $[0,1]$ plus labor-share automation, augmentation, mitigation, and partial-adjustment parameters & Indicative provisional grids $\delta_s,\psi_s,\zeta_s\in\{0.25,0.50,1.00\}$ and $\lambda_{LS}\in\{0.25,0.50,1.00\}$; enforce $\zeta_sRST\leq\delta_s$ outside augmentation-dominant robustness; finalized in the baseline calibration registry (replication package, appendix\_i files) & Low/central/high grid for scenario parameters; $RST\in[0,1]$ from declared proxy; bounded PERT only after pilot supports are declared & \citet{acemoglu2018race,acemoglu2019automation}; active labour-market policy spending, training participation (ILOSTAT), or residualized human-capital component of AIPI; calibrated scenario discipline \\
$MFC^{gross}_{i,s,r,t}$ & Gross marginal fiscal capture of AI-induced output before administrative costs, transition costs, and leakage & Construct from channel model or draw from historical gross-capture distribution, but not both in the same run & Channel-mode derived object or historical-mode stochastic object; report mode explicitly & UNU-WIDER GRD; \citet{brollo2024fiscal}; IMF fiscal-space guidance \\
$MFC^{eff}_{i,p,s,r,t}$ & Derived reporting object after subtracting $ac$, $tr$, and $leak$ from fiscal space & Compute only as $MFC^{eff}=fs^{eff}/g^{AI}$ when $g^{AI}>0$; not drawn as an independent primitive & Report for interpretation; do not use for gross historical plausibility checks & Derived from $MFC^{gross}$, $g^{AI}$, $ac$, $tr$, and $leak$ \\
$\tau^{L,fisc,eff},\tau^{K,fisc,eff},\tau^{C,eff},\tau^{AI}$; $\tau^{L,disp},\tau^{K,disp}$ & Fiscal-capture rates for revenue and disposable-income rates for induced consumption & Infer fiscal rates from revenue composition where possible; disposable rates may coincide only in simplified calibration & Grid by regime; cap AI-rent tax at plausible maximum $\bar\tau^{AI}$ & GRD revenue composition; \citet{brollo2024fiscal} \\
$\chi^{Kbase},\chi^{dom},\psi^{shift}$; $\psi^{K,rate}$ & Capital-base chain parameters, each in $[0,1]$; $\psi^{K,rate}$ is the rate-side shifting parameter, active only under the labelled rate-side convention with $\psi^{shift}=0$ & Construct the taxable domestic post-shifting capital base without double counting & Bounded PERT on low/central/high supports & National accounts and profit-shifting literature \\
$mpc^W,mpc^{\Pi},mpc^R,mpc^{ben},m^M,exempt^C,\theta^{R,dom},\tau^{R,disp}$ & Induced-consumption shares and rates, each in $[0,1]$; $mpc^{ben}$ active only in the labelled first-round recycling variant & Calibrate disposable-income pass-through, domestic routing, imports, and exemptions consistently & Bounded PERT or declared grid & Household surveys, national accounts, and VAT rules \\
$\lambda^{AI},\chi^{AI,dom},\psi^{shift,tech},\chi^{AI,compliance},\chi^{AI,enforce}$ & AI-rent share and capturability components in $[0,1]$ & Zero AI-rent channel in the status quo through $\tau^{AI}=0$ unless such an instrument already exists; bounded shares in positive scenarios & Bounded scenario grid or PERT & Fiscal law, international taxation, and Appendix \ref{app:fiscal_conversion_module} \\
$I^L_{i,t}, I^F_{i,t}, I^C_{i,t}$ & Labour, firm, and consumption-base informality & Use source-specific measures where available; do not collapse all frictions into one informal-economy scalar unless data force it & Fixed observed value; marginal informality $I^{j,marg}$ preferred when sectoral AI-shock weights exist, aggregate $I^j$ as declared fallback & ILOSTAT; \citet{worldbank2021informality,medina2018shadow} \\
$\rho_L,\rho_K,\rho_C,\rho_E,\rho_\psi$ & Elasticities for capturability, exemptions, and rate-side profit shifting & No preferred point value without estimation; $\rho_E$ and $\rho_\psi$ active only in their labelled conventions & Grid $\{0.5,1,1.5,2\}$; report sensitivity and Spearman/Sobol diagnostics & Calibrated structural parameters \\
$\epsilon^{ero,j}$ & Behavioral taxable-base erosion elasticities by channel & Zero only in the mechanical upper-bound reading of reform regimes; mandatory companion for $r\geq1$ results, with $\epsilon^{ero,j}\tau^j<1$ & Declared grid or bounded PERT; recorded in the calibration audit trail & Tax avoidance and evasion elasticity literature; Appendix \ref{app:fiscal_conversion_module} \\
$ac^{net,fix},tr^{ann},leak^{fix}$ & Fixed administrative, transition, and residual leakage costs as GDP-normalized ratios & Subtract from fiscal space after the AI shock is converted into revenue & Bounded PERT supports; triangular only as robustness & ASPIRE expenditure concepts; IMF fiscal policy; GRD \\
$ac^{ME},tr^{ME},leak^{ME}$ & Marginal-equivalent deductions per unit of AI-induced output & Enter only through $\widetilde{MFC}^{gross}=MFC^{gross}-ac^{ME}-tr^{ME}-leak^{ME}$ & Bounded PERT or zero in baseline if costs are fixed & IMF fiscal policy; implementation stress tests \\
$\Omega^{res},s^{PB}$ & Residual stress penalty and primary-balance correction & $\Omega^{res}=0$ in the baseline; derive $s^{PB}$ from baseline and debt-stabilizing primary balances & Positive stress grid for $\Omega^{res}$; WEO/DSA scenarios for $s^{PB}$ & IMF WEO and debt-sustainability analysis \\
$\xi$, $\bar\pi$ & Prudential buffer and probability label & $\xi\in\{0,0.05,0.10,0.25,0.50\}$ and $\bar\pi\in\{0.50,0.75,0.90\}$ & Fixed reporting grid, not estimated parameter & Robust Decision Making logic \citep{lempert2003shaping,rand2013rdm} \\
\bottomrule
\end{longtable}
\normalsize

\subsection{Data construction rules}

Observed variables should be constructed as follows:
\begin{itemize}
    \item Fiscal capacity: tax revenue, non-tax revenue, grants, social contributions, and revenue composition from the UNU-WIDER Government Revenue Dataset \citep{unu2025grd}; use central/general government consistently.
    \item Informality: informal employment as a share of total employment from ILOSTAT where available, supplemented by shadow-economy or informality studies only when national labour-force measures are missing \citep{ilostat2026informal,medina2018shadow,worldbank2021informality}.
    \item Social protection: coverage, adequacy, expenditure, beneficiary information, and program categories from ASPIRE where available \citep{worldbank2026aspire}.
    \item Poverty lines and poverty gaps: World Bank Poverty and Inequality Platform or national poverty lines, with PPP and local-currency concepts kept separate \citep{worldbank2026pip}.
    \item AI preparedness: IMF AI Preparedness Index and its four components; do not reinterpret the index as invention leadership \citep{cazzaniga2024genai}.
    \item AI exposure: occupation- or sector-level exposure harmonized to national employment or value-added shares. Productive exposure and displacement exposure must be reported as separate variables.
    \item Digital and skills gap: construct $Gap_{i,t}$ from ITU and World Bank digital indicators excluded from AIPI, and document the normalization to $[0,1]$.
    \item Primary-balance correction: obtain $PB^0_{i,t}$ and $PB^{stab}_{i,t}$ for $s^{PB}_{i,t}$ from the IMF WEO and debt-sustainability analysis.
    \item External translation frontier: document the benchmark economies and year used to construct $S^{frontier}_{s,t}$, which remains fixed within each threshold inversion, and state that the benchmark set matches the population underlying the scenario shock $\phi_s$ (population consistency).
\end{itemize}

\subsection{Maintained economic assumptions and bias directions}

The accounting frontier rests on maintained assumptions stated throughout the paper. They are consolidated here with their bias direction: conservative assumptions understate feasibility, optimistic ones overstate it.

\begin{longtable}{p{0.43\textwidth}p{0.22\textwidth}p{0.25\textwidth}}
\caption{Maintained economic assumptions and bias directions}\label{tab:maintained_assumptions_bias}\\
\toprule
Assumption & Stated in & Direction \\
\midrule
\endfirsthead
\toprule
Assumption & Stated in & Direction \\
\midrule
\endhead
One-directional chain; no feedback from protection to productivity & Section \ref{sec:conceptual} & Ambiguous \\
Endpoint level gain persists as long as the policy; composition may degrade through rent dissipation & Section \ref{subsec:phi_ladder_main} & Optimistic for long-run capture \\
Non-market and zero-price AI gains outside the fiscal base; framework measures taxable income, not welfare & Sections \ref{subsec:phi_ladder_main} and \ref{sec:claims} & Conservative \\
\(\kappa^{TFP\to Y}\leq1\) excludes induced capital deepening & Section \ref{app:relocated_main} & Conservative \\
No general equilibrium: no wage spillovers, first-round consumption only & Sections \ref{sec:conceptual} and \ref{sec:fiscal_conversion} & Conservative \\
\(S^{frontier}\) benchmark matches the population of \(\phi_s\) & Section \ref{subsec:ndc_translation_main} & Conservative wedge if mismatched \\
No endogenous imitation diffusion; adoption paths are scenario paths & Appendix \ref{app:adoption_module} & Conservative in the medium term \\
Informality fixed in baseline; formalization and informalization only as labelled scenarios, reported jointly & Appendix \ref{app:informality_module} & Ambiguous \\
One adoption stock gates automation and augmentation; \(E^{auto}\) anchored to full-automation risk is a lower bound on displacement & Section \ref{sec:labor_share} & Optimistic for the labor tax base \\
\(\chi^{Kbase}\) calibrated on average non-labor-residual composition and applied to the AI increment & Appendix \ref{app:labor_share_module} & Conservative \\
LAC anchoring uses productivity-enhancing exposure only; automation-driven output gains excluded from \(E^{prod}\) & Section \ref{subsec:prod_disp_exposure_main} & Conservative \\
Statutory-accrual incidence; no incidence shifting & Section \ref{sec:fiscal_conversion} & Ambiguous \\
Reform-regime gains are mechanical upper bounds; behavioral-erosion companion mandatory for \(r\geq1\) & Section \ref{sec:fiscal_conversion} & Optimistic for \(r\geq1\) \\
Historical capture prior is composition-weighted and policy-inclusive (buoyancy, not elasticity) & Section \ref{sec:historical_plausibility} & Ambiguous; generous for $r=0$ plausibility labels \\
Recipient behavior fixed; GMI priced at ideal targeting absent \(\theta^{target}\); no transfer recycling & Section \ref{sec:policy_costs} & GMI-optimistic; recycling omission conservative \\
Feasibility is an endpoint condition; crossing year is not a launch year & Section \ref{sec:thresholds} & Naive launch reading optimistic \\
\(PB^{stab}\) static at baseline growth & Section \ref{sec:thresholds} & Conservative \\
\bottomrule
\end{longtable}

\section{Additional derivations and protocols relocated from the main text}\label{app:relocated_main}

\subsection{Unit bridges and AI-shock objects}
\begingroup
\let\subsection\subsubsection
% S3-SPAN-M2-BEGIN
Let $\phi^{raw}_{s,u}$ denote the raw AI shock reported by a source under unit $u\in\{TFP,LP,Y\}$, where $TFP$ is total factor productivity, $LP$ is labour productivity, and $Y$ is GDP-equivalent output. The model-ready real GDP-equivalent shock is obtained only after an explicit bridge:
\begin{equation}
    \phi^Y_s=\kappa^{u\to Y}_s\phi^{raw}_{s,u},\qquad 0<\kappa^{u\to Y}_s\leq 1\ \text{(baseline convention)}.
\end{equation}
For $u=Y$, the bridge is the identity, $\kappa^{Y\to Y}=1$. For $u=LP$, $\kappa^{LP\to Y}\leq1$ reflects that GDP gains cannot exceed labour-productivity gains unless employment rises. For $u=TFP$, standard growth accounting implies amplification through capital deepening ($\Delta\ln Y=\Delta\ln TFP/(1-\alpha_K)>\Delta\ln TFP$); the baseline restriction $\kappa^{TFP\to Y}\leq1$ is therefore a deliberate conservative convention that excludes induced capital accumulation from the fiscal-base shock. A capital-deepening robustness bridge $\kappa^{TFP\to Y}\in[1,1/(1-\alpha_K)]$ may be reported separately and must be labelled as such.
The nominal fiscal-base-compatible shock is then
\begin{equation}
    \phi^{Y,nom}_s=(1+\phi^Y_s)(1+\pi^{AI,Y}_s)-1,
\end{equation}
or, for small rates, $\phi^{Y,nom}_s\approx \phi^Y_s+\pi^{AI,Y}_s$. The exact expression is multiplicative, not log-additive, and $\pi^{AI,Y}_s$ may be negative. The nominal bridge is admissible only when fiscal-capture anchors and policy costs are nominal and share the endpoint horizon. For price-indexed benefits, poverty lines, or eligibility costs, the baseline instead pairs a real GDP-equivalent shock with real capture anchors or projects costs to that nominal horizon; price effects alone are never real fiscal capacity. Non-market or zero-price AI output lies outside GDP and the fiscal base, distinct from informality, biasing the frontier conservatively. Although source-unit supports are positive, draws with $\phi^{Y,nom}_s\leq0$ remain economic failure draws in all-draw probabilities and are excluded only from inversions conditional on $\phi^{Y,nom}_s>0$ (Section \ref{subsec:impossibility_score}). A persistently negative aggregate effect would reduce $fs^{eff}$ further, so the reported frontiers are upper bounds on that margin.

\begin{table}[H]
\centering
\caption{Mandatory AI-shock objects and admissible use}
\label{tab:ai_shock_objects_v17}
\small
\begin{tabular}{p{0.25\textwidth}p{0.48\textwidth}p{0.18\textwidth}}
\toprule
Object & Meaning & Admissible use \\
\midrule
$\phi^{raw}_{s,u}$ & Literature shock in source unit $u\in\{TFP,LP,Y\}$ & Input only; never used directly in fiscal-space equations \\
$\phi^{Y,nom}_{s,t}$ & Annual nominal GDP-equivalent, fiscal-base-compatible shock after unit and nominal bridging & Annual increment to be accumulated or used in a labelled static approximation \\
$g^{AI,level}_{i,s,H}$ & AI-induced level gain relative to the no-AI counterfactual over horizon $H$ & Object compared with recurring annual policy costs \\
\bottomrule
\end{tabular}
\end{table}
The preferred stochastic assignment applies bounded PERT to the source-unit shock:
\begin{equation}
    \phi^{raw}_{s,u}\sim \mathrm{PERT}(a_s,m_s,b_s).
\end{equation}
Conversion then follows the bridges above; $a_s$ and $b_s$ are the table's source-unit bounds, and $m_s$ is their midpoint unless a source-specific central estimate is available.

% S3-SPAN-M2-END
\endgroup

\subsection{Logistic adoption and diffusion module}
\begingroup
\let\subsection\subsubsection
% S3-SPAN-M3-BEGIN
\subsection{Calibrable logistic adoption module}\label{subsec:logistic_adoption_main}

The framework distinguishes technical access, task-level use, and productive adoption integrated into output processes; micro productivity gains need not become macro gains.

The baseline use equation is
\begin{equation}\label{eq:q_use_logistic_main_v10}
    q^{use}_{i,s,t}=\bar q_{i,t}\Lambda(\ell^q_{i,s,t}),
    \qquad
    \Lambda(x)=\frac{1}{1+e^{-x}}.
\end{equation}
To avoid double counting productive exposure, the baseline adoption index does not include $E^{prod}$ in the logistic adoption equation. Productive exposure enters the translation score directly, while adoption is restricted to preparedness and frictions:
\begin{equation}\label{eq:eta_noE_main_v10}
    \ell^q_{i,s,t}=\mu_{i,s,t}+\nu_A A_{i,t}-\nu_I I_{i,t}-\nu_G Gap_{i,t}.
\end{equation}
Here $A$ is AI preparedness, $I$ is the relevant adoption-side informality or organizational-friction index, and $Gap$ captures digital-access, skills, or infrastructure gaps not already embedded in $A$. In the baseline, $Gap$ is constructed only from indicators excluded from the AIPI preparedness index. If overlapping indicators are unavoidable, the implementation uses a residualized gap, $Gap^{\perp}_{i,t}=Gap_{i,t}-\operatorname{Proj}_{A}(Gap_{i,t})$, or moves the overlapping gap measure to robustness. A robustness specification may include $E^{prod}$ in $\ell^q$, but then exposure must not also enter the translation score as an independent multiplicative factor.

The country adoption ceiling is
\begin{equation}\label{eq:qbar_main_v10}
    \bar q_{i,t}=\max\left\{0,\min\left[1,1-\omega_I I_{i,t}-\omega_G Gap_{i,t}\right]\right\}.
\end{equation}
Informality and $Gap$ may enter both short-run diffusion and long-run adoption capacity only if their empirical content is partitioned across the diffusion and ceiling equations or one margin is residualized against the other. This is an anchored calibration, not an identified cross-country forecasting equation; unless a time trend or scenario path is explicitly calibrated, the anchored intercept is held fixed after $t_0$, $\mu_{i,s,t}=\mu_{i,s,t_0}$ for $t>t_0$. The logit inversion of the intercept, interior projection of the target, initial condition, and support restrictions are specified in Online Appendix B.
Productive adoption follows a lagged diffusion law with an explicit ceiling projection:
\begin{equation}\label{eq:q_prod_diffusion_main_v10}
    q^{prod}_{i,s,t}=\min\left\{\bar q_{i,t},
    (1-\lambda_q)q^{prod}_{i,s,t-1}+\lambda_q q^{use}_{i,s,t}\right\},
    \qquad \lambda_q\in\{0.15,0.30,0.50\}\quad\text{in baseline, with }\lambda_q=1\text{ only as an instant-adoption stress test}.
\end{equation}
The projection is needed because $q^{prod}$ is stock-like; if the ceiling $\bar q_{i,t}$ falls after a shock to infrastructure, informality, or access, the unprojected stock could otherwise exceed the current feasible ceiling.
% S3-SPAN-M3-END
\endgroup

\subsection{Capital-base and domestic exposure chain}
\begingroup
\let\subsection\subsubsection
% S3-SPAN-M5-BEGIN
The accounting chain is
\begin{equation}\label{eq:capital_base_chain_v17}
\begin{split}
    \Delta NLS^{AI}_{i,s,t}
    &\rightarrow \Delta\Pi^{AI,gross}_{i,s,t}
    =\chi^{Kbase}_{i,t}\Delta NLS^{AI}_{i,s,t},\\
    \Delta\Pi^{AI,gross}_{i,s,t}
    &\rightarrow \Delta\Pi^{AI,dom}_{i,s,t}
    =\chi^{dom}_{i,s,t}(1-\psi^{shift}_{i,s,t})\Delta\Pi^{AI,gross}_{i,s,t},\\
    \Delta\Pi^{AI,dom}_{i,s,t}
    &\rightarrow \Delta\Pi^{AI,disp}_{i,s,t}
    =(1-\tau^{K,disp}_{i,s,r,t})\Delta\Pi^{AI,dom}_{i,s,t}.
\end{split}
\end{equation}
The capital fiscal-capture rate $\tau^{K,fisc,eff}_{i,s,r,t}$ and disposable-income rate $\tau^{K,disp}_{i,s,r,t}$ may differ because compliance, corporate taxation, withholding, and related institutions need not map one-to-one into household income. Consumption is an induced base from disposable income and import leakages, not a third functional output component. Thus the fiscal weights are tax-base exposure ratios, not exhaustive shares of $\Delta Y^{AI}$, and need not sum to one.

In particular, the capital/profit exposure weight is constructed on the domestic, profit-shifting-adjusted taxable-profit base:
\begin{equation}\label{eq:omega_k_dom_base_v25}
\begin{split}
    \omega^K_{i,s,t}
    &=\frac{\Delta\Pi^{AI,dom}_{i,s,t}/Y^0_{i,t}}{g^{AI}_{i,s,t}} \\
    &=\chi^{dom}_{i,s,t}(1-\psi^{shift}_{i,s,t})\chi^{Kbase}_{i,t}
      \frac{\Delta NLS^{AI}_{i,s,t}/Y^0_{i,t}}{g^{AI}_{i,s,t}},
\end{split}
\end{equation}
whenever $g^{AI}_{i,s,t}>0$. If $g^{AI}_{i,s,t}=0$, the draw is treated as an economic non-crossing draw and ratio-based exposure weights are not computed. Equation \eqref{eq:omega_k_dom_base_v25} is the baseline convention: foreign ownership and ordinary profit shifting live in the base-side construction of $\Delta\Pi^{AI,dom}$ and are not multiplied again in $\tau^{K,fisc,eff}$ or $\chi^K$. A rate-side profit-shifting robustness convention is allowed only if $\psi^{shift}$ is set to zero in Equation \eqref{eq:capital_base_chain_v17}.

% S3-SPAN-M5-END
\endgroup

\subsection{Consumption, fiscal anchors, and calibration}
\begingroup
\let\subsection\subsubsection
% S3-SPAN-M5c-BEGIN
\subsection{Consumption as an induced tax base}

Consumption is an induced tax base rather than a third functional component of output. The disposable-income rates $\tau^{j,disp}$ are not fiscal-capture rates, and import content and consumption exemptions enter exactly once, on either the base side or the rate side. The consumption channel is first-round only: induced consumption is not recycled into further income and consumption rounds. This truncation is consistent with the absence of general-equilibrium multipliers elsewhere in the framework and biases the frontier conservatively. Online Appendix D gives the full construction.

\subsection{Observed fiscal anchors}

The preferred fiscal data source for cross-country revenue anchors is the UNU-WIDER Government Revenue Dataset, which provides central and general government revenue series and, where possible, figures with and without natural-resource revenues \citep{unu2025grd}. OECD Revenue Statistics in Latin America and the Caribbean is used as a harmonized tax-to-GDP cross-check for Peru, Chile, Colombia, and Mexico \citep{oecd2025revenue}. The baseline empirical implementation records three fiscal anchors:
\begin{equation}
    \theta^{tax}_{i,t}=\frac{TaxRev_{i,t}}{Y_{i,t}},\qquad
    \theta^{rev}_{i,t}=\frac{TotalRev_{i,t}}{Y_{i,t}},\qquad
    \theta^{nonres}_{i,t}=\frac{Rev^{nonres}_{i,t}}{Y_{i,t}}.
\end{equation}
The non-resource anchor is preferred for countries where commodity cycles can make fiscal capacity look temporarily high. These anchors discipline $MFC^{gross}$; they are not automatically equal to $MFC^{eff}$.

\subsection{Country-specific calibration rule}

For each country $i$, fiscal regime $r$, and scenario $s$, fiscal-conversion draws are constructed in two steps. First draw or compute gross marginal capture. For a positive-capture conditional feasibility benchmark,
\begin{equation}
    MFC^{gross}_{i,s,r,t}
    \sim \text{PERT}\left(P10(MFC^{hist,gross,+}_i),P50(MFC^{hist,gross,+}_i),P90(MFC^{hist,gross,+}_i)\right),
\end{equation}
with country-regime bounds. A tighter $\text{PERT}(P25,P50,P75)$ is a low-dispersion robustness variant, not the baseline, because quartile bounds truncate support more than decile bounds. Unconditional assessment uses the two-component mixture in Section \ref{sec:historical_plausibility}. Unit matching follows the canonical nominal/real rule in Section \ref{subsec:phi_ladder_main}; in particular, the baseline never combines real/TFP-equivalent $g^{AI}$ with nominal historical $MFC^{gross}$.

Then compute $fs^{eff}=\widetilde{MFC}^{gross}g^{AI,level}-ac^{net,fix}-tr^{ann}-leak^{fix}$ and derive $MFC^{eff}=fs^{eff}/g^{AI,level}$ only when needed for reporting and $g^{AI,level}>0$. This ensures that the historical plausibility check compares gross objects with gross objects, while effective feasibility uses net fiscal space.


% S3-SPAN-M5c-END
\endgroup

\subsection{Historical capture mixture and cleaning rules}
\begingroup
\let\subsection\subsubsection
% S3-SPAN-M6-BEGIN
Because $\Delta R$ may be negative even when nominal GDP growth is positive, the implementation reports both the raw historical distribution and a non-negative-capture distribution:
\begin{equation}\label{eq:mfc_hist_positive_v10}
    H_i^+=\{MFC^{hist,gross}_{i,t}:MFC^{hist,gross}_{i,t}\geq0\}.
\end{equation}
Negative-capture years remain reported as fiscal-stress years. The positive-capture calibration $MFC^{gross}\sim MFC^{hist,gross}\mid MFC^{hist,gross}\geq0$ is explicitly conditional and labelled as a positive-capture scenario, not unconditional baseline capacity.
The empirical table must also report
\begin{equation}
    p_i^- = \Pr\left(MFC^{hist,gross}_{i,t}<0\right)
\end{equation}
inside the historical window. When the goal is unconditional uncertainty propagation, the preferred reduced-form draw is a two-component mixture,
\begin{equation}
    MFC^{gross}\sim
    \begin{cases}
    MFC^{hist,gross,-}, & \text{with probability }p_i^-,\\
    MFC^{hist,gross,+}, & \text{with probability }1-p_i^- ,
    \end{cases}
\end{equation}
Here $MFC^{hist,gross,+}$ and $MFC^{hist,gross,-}$ are the positive and fiscal-stress supports. PERT draws from $H_i^+$ are positive-capture conditional and may overstate unconditional capacity in crisis-prone economies. Nor does this imply $MFC^{eff}\geq0$: net fiscal space can be negative after costs and leakage, and negative $fs^{eff}$ draws are economic failures, not invalid draws.
The mixture is unconditional over the sign of capture but conditional on positive nominal GDP growth (cleaning rule 2), the relevant class for a positive AI-induced increment. Zero or negative increments remain economic non-crossings or failure draws without a capture ratio.

The cleaning rules are:
\begin{enumerate}
    \item Use one government level consistently within each country.
    \item Exclude years with non-positive nominal GDP growth when computing increment ratios.
    \item Report raw and positive-capture distributions separately.
    \item Report results with and without pandemic years 2020--2021.
    \item Report results with and without resource revenue where GRD provides the distinction.
    \item Winsorize raw ratios at $P1/P99$ or report robust medians to avoid denominator artifacts.
\end{enumerate}
% S3-SPAN-M6-END
\endgroup

\subsection{Financeable policy-cost envelopes}
\begingroup
\let\subsection\subsubsection
% S3-SPAN-M7-BEGIN
The maximum financeable gross policy cost is reported as two separate non-negative policy envelopes. The marginal accounting envelope is
\begin{equation}\label{eq:cmax_main_v17}
    c^{gross,max,+}_{i,p,s,r,t}(\xi)=\max\left\{0,\frac{fs^{eff}_{i,p,s,r,t}}{1+\xi}\right\}.
\end{equation}
The stronger debt-consistent envelope is
\begin{equation}\label{eq:cmax_dc_main_v22}
    c^{gross,max,+,DC}_{i,p,s,r,t}(\xi)=\max\left\{0,\frac{fs^{eff}_{i,p,s,r,t}-s^{PB}_{i,t}}{1+\xi}\right\}.
\end{equation}
If the relevant numerator is negative, no positive policy cost is financeable under that scenario and regime.

% S3-SPAN-M7-END
\endgroup

\subsection{Parameter hierarchy and assignment details}
\begingroup
\let\subsection\subsubsection
% S3-SPAN-M8-BEGIN
\begin{longtable}{p{0.18\textwidth}p{0.30\textwidth}p{0.34\textwidth}}
\caption{Parameter hierarchy and assignment rules}\label{tab:parameter_hierarchy}\\
\toprule
Class & Examples & Assignment rule \\
\midrule
\endfirsthead
\toprule
Class & Examples & Assignment rule \\
\midrule
\endhead
Observed country data & GDP, population, eligible population, tax revenue/GDP, primary balance, informality, poverty lines, social protection expenditure & Use public data directly after unit harmonization; no scenario assignment unless the data are missing or stale \\
Constructed indices & AI preparedness, productive exposure, digital gap, fiscal capturability, policy cost net of existing transfers & Report raw indicators, normalization, weights, and aggregation formula; no hidden transformations \\
Literature-disciplined scenarios & Base AI shock $\phi_s$, exposure ranges, micro productivity relevance, diffusion speeds & Anchor low/central/high values in published studies; report the source-specific unit before conversion \\
Calibrated structural parameters & Translation elasticities $\alpha,\beta,\gamma$, informality elasticities $\rho_L,\rho_K,\rho_C$, leakage, AI-rent share, rent-capture rate & Use grids and Monte Carlo; no single preferred value unless estimated from data \\
Stress-test parameters & Disruptive AI shock, instant adoption, zero leakage, high rent taxation, high formalization effect & Use only to identify frontier requirements; do not present as baseline forecasts \\
\bottomrule
\end{longtable}
\subsection{Distributional assignment rules}

When a parameter has a lower value $a$, a modal or central value $m$, and an upper value $b$, the baseline uncertainty distribution is a bounded PERT distribution. If $a=b$, the parameter is treated as fixed at that value and no beta draw is attempted:
\begin{equation}
    X=a+(b-a)Z,
    \qquad
    Z\sim Beta\left(1+4\frac{m-a}{b-a},\ 1+4\frac{b-m}{b-a}\right).
\end{equation}
With only a defensible support interval, use a uniform distribution and state that no central value is imposed. A lognormal is allowed only for strictly positive multiplicative parameters, such as leakage or administrative cost, after reporting its median and 90 percent interval. Bounded shares are truncated to their natural support.

Both independence and scenario-consistent dependence specifications are reported. Without an empirical correlation matrix, correlations are scenario assumptions, not estimates; Online Appendix G gives the dependence structure.

\subsection{Data construction rules}

Observed variables are constructed from declared primary sources: GRD, ILOSTAT, ASPIRE, PIP or national poverty lines, AIPI, ITU and World Bank indicators excluded from AIPI, and WEO/DSA. Online Appendix H gives the complete construction rules. Executed choices are recorded in the calibration registry and the Results data section.

\subsection{Calibration audit trail}

Each empirical table should be accompanied by a machine-readable calibration file with the following fields:
\begin{equation}
\begin{gathered}
    \{\text{parameter, country, year, raw source, raw value, transformation, unit,}\}\\
    \{\text{support, distribution, dependence block, scenario class}\}.
\end{gathered}
\end{equation}
This audit trail is not optional. It is what makes the framework reusable when future AI productivity estimates, adoption measures, or fiscal datasets change.

% S3-SPAN-M8-END
\endgroup

\subsection{Diagnostic protocols}
\begingroup
\let\subsection\subsubsection
% S3-SPAN-M9b-BEGIN
\subsection{Temporal placebo}

The placebo protocol has two layers. The first is a mechanical accounting-null check. In pre-GenAI years, the GenAI-specific productivity shock is set to zero,
\begin{equation}
    \phi^{GenAI}_{s,t}=0\qquad \text{for accounting-null placebo years.}
\end{equation}
With the zero-adoption gate active, this implies $T^{GenAI}=0$ and $g^{AI,level}=0$ for the GenAI channel. This null check verifies that the code does not manufacture AI fiscal space mechanically, but it is not interpreted as an informative historical placebo because $0$ minus positive deployment costs is arithmetically negative.

The second layer is the informative pre-GenAI ICT placebo. For years such as 2010--2018, the implementation feeds the translation block with historical digitalization or ICT-adoption proxies rather than GenAI adoption:
\begin{equation}
    S^{ICT}_{i,t}=\mathbf{1}\{q^{ICT}_{i,t}>0,\ E^{ICT}_{i,t}>0\}(\tilde E^{ICT}_{i,t})^{\beta_{ICT}}(\tilde q^{ICT,num}_{i,t})^{\gamma_{ICT}},
    \qquad
    T^{ICT}_{i,t}=\min\left\{1,\frac{S^{ICT}_{i,t}}{S^{ICT,frontier}_{t}}\right\}.
\end{equation}
The placebo shock is then
\begin{equation}
    g^{ICT,placebo}_{i,t}=\phi^{ICT,placebo}_{t}T^{ICT}_{i,t},
\end{equation}
where $\phi^{ICT,placebo}$ is a documented growth-accounting ICT/digitalization analogue, not a GenAI shock, with comparability defaults $\beta_{ICT}=\beta$ and $\gamma_{ICT}=\gamma$. A strong flag marks broad pre-GenAI placebo feasibility of AI-funded UBI/PEN/GMI under the same thresholds:
\begin{equation}
    PlaceboFail^{ICT}_{i,p,t}=\mathbf{1}\{V^{gross,ICT}_{i,p,t}\geq1\ \text{for broad policy cells in pre-GenAI years}\}.
\end{equation}
The test asks whether ordinary digitalization, growth, or buoyancy is being relabelled as AI fiscal space; it is diagnostic, not causal. Its known-outcome benchmark is that digitalization in 2010--2018 financed universal social protection nowhere in the region, so broad feasibility contradicts history, not merely a prior. Validity requires the full era-consistent pipeline---historical capture distributions, policy costs, and informality---not translation alone.

\subsection{Sectoral negative controls}

For sectors with very low AI exposure, the translation score should be low. Let $\mathcal{K}^{NC}$ be a set of low-exposure sectors. A negative-control diagnostic reports
\begin{equation}
    T^{NC}_{i,s,t}=\sum_{k\in\mathcal{K}^{NC}}\omega^{NC}_{i,k,t}T_{i,k,s,t}.
\end{equation}
The weights $\omega^{NC}_{i,k,t}$ are employment or value-added shares within the negative-control set and satisfy $\sum_{k\in\mathcal{K}^{NC}}\omega^{NC}_{i,k,t}=1$.
The expectation is not that $T^{NC}=0$, but that it remains substantially below the country aggregate in scenarios where the aggregate is driven by high-exposure sectors. If negative-control sectors generate high translation scores, the exposure/adoption construction is too permissive. These placebo and negative-control diagnostics are diagnostic stress tests, not substitutes for a causal treatment-effect design \citep{bertrand2004trust,athey2017state,lipsitch2010negative}.

\subsection{Leave-one-source-out and channel-ablation protocol}

The replication package reports leave-one-source-out and channel-ablation checks:
\begin{itemize}
    \item AI-productivity source: Acemoglu-only, OECD-only, and high/stress-only exclusions.
    \item Exposure source: IMF/OECD exposure only, World Bank--ILO exposure only, and country-sector exposure only where available.
    \item Fiscal channel: no AI-rent tax, no formalization effect, no labor augmentation, high import leakage, and no base-broadening reform.
    \item Cost convention: gross-cost baseline, net-cost robustness, and debt-consistent requirement.
\end{itemize}
If a positive feasibility result depends on a single source or a single fiscal channel, the result is reported as source-dependent rather than robust.


% S3-SPAN-M9b-END
\endgroup

\subsection{Monte Carlo propagation and sensitivity protocols}
\begingroup
\let\subsection\subsubsection
% S3-SPAN-M10-BEGIN
\subsection{Monte Carlo uncertainty propagation}\label{subsec:mc_main}

Let $\Theta^{(m)}$ denote draw $m$ of the uncertain-input vector. The full vector must state which variables are observed and fixed and which are stochastic. The implementation must choose one fiscal-conversion mode and one adoption mode for a given run.

Observed-adoption mode draws $q^{prod}$ from adoption or measurement-error distributions and treats $A,I,Gap$ as controls. Anchored simulated-adoption mode takes $A,I,Gap,\lambda_q,\nu_A,\nu_I,\nu_G,q^{use,target}_{i,t_0}$ as primitives; logit-inverted $\mu_{i,s,t_0}$ is not an independently identified structural coefficient. Structural-intercept robustness instead draws a common or hierarchical intercept without target inversion. Simulated-adoption modes derive $q^{use}$ and $q^{prod}$ endogenously and never draw $q^{prod}$ independently.

Channel mode derives gross capture from tax bases rather than drawing it. Threshold runs declare frozen versus endogenous weights; only the former use closed forms, while the latter solve the channel-revenue function in Online Appendix D. Online Appendix G gives both mode vectors and the complete AI-rent primitives $\tau^{AI},\chi^{AI,dom},\psi^{shift,tech},\chi^{AI,compliance},\chi^{AI,enforce}$, with derived $\tau^{AI,eff}=\tau^{AI}\chi^{AI}$; historical reduced-form mode instead draws gross capture.
Capturability elasticities $\rho_j$ are outside $\Theta_{hist}$. Labelled alternatives add their parameters to the audit trail: $\psi^{K,rate}$ for rate-side shifting and $\epsilon^{ero,j}$ for behavioral erosion.
Historical-mode channel rates may be descriptive but do not also construct $MFC^{gross}$. The run declares observed or simulated adoption; only observed-adoption mode treats $q^{prod}$ as primitive. Fixed GDP-share costs and marginal-equivalent deductions each occupy one margin. Nominal $AC$, $TR$, and $Leak$ are derived from GDP, never drawn independently alongside normalized objects.

For each draw, effective fiscal space is computed directly from net marginal gross capture and fixed costs; when $g^{AI,level,(m)}>0$, the per-unit object $MFC^{eff,(m)}=fs^{eff,(m)}/g^{AI,level,(m)}$ is derived rather than drawn. The complete draw-level specification is in Online Appendix G.
A draw with $fs^{eff,(m)}<0$ is not invalid. It is an economically meaningful failure draw: administrative costs, transition costs, and leakage exceed gross AI-generated fiscal space. Such draws remain in the all-draw feasibility probability and lower the crossing probability.
Arithmetic, computability, support, numerical-floor, and failure semantics follow the canonical conservative reporting chain in Section \ref{subsec:impossibility_score}.
The gross feasibility ratio is $V^{gross,(m)}=fs^{eff,(m)}/c^{gross,(m)}$. The net ratio $V^{net,(m)}=fs^{eff,(m)}/c^{net,(m)}$ is computed only for draws with $c^{net,(m)}>0$.
If $c^{net,(m)}=0$, the draw is classified as already covered by existing spending on the net-cost margin. $V^{gross}$ remains the baseline reporting object.

\subsection{Monte Carlo output layers}

For each country-policy-scenario-regime cell, the percentile and conditional-probability layers are those specified in Online Appendix G; the canonical all-draw crossing grid is Equation \eqref{eq:mc_prob_main_v10}.
If the denominator is zero, the corresponding conditional probability is reported as undefined/NA and the relevant valid-share, $\pi^{basic}$ or $\pi^{support}$, is reported as zero. This distinguishes ``no valid support exists'' from ``valid support exists but no draw crosses the threshold.'' The debt-consistent version uses Equation \eqref{eq:debt_requirement_main_v10}.

\subsection{Distributional choices and dependence}

The baseline distribution for low-central-high calibrated scalar parameters is bounded PERT. Triangular and truncated-uniform draws are robustness checks, not the baseline convention. Bounded share variables may also be drawn from beta distributions when moments are available. Dependence between preparedness, adoption, informality, and fiscal capacity may be implemented through correlated latent shocks. If dependence is imposed, the paper reports Spearman and tornado diagnostics as the baseline sensitivity outputs and reserves classical Sobol indices for an independent-input diagnostic run or a dependence-aware sensitivity method; Online Appendix G gives the scenario-consistent dependence structure.

\subsection{Global sensitivity and stress tests}

Report Spearman rank correlations, one-at-a-time tornado ranges, and grouped block sensitivity for AI productivity, adoption, fiscal capture, informality, leakage, and cost \citep{mckay1979lhs,morris1991factorial,saltelli2008global,sobol2001global}. Classical Sobol indices are reserved for independent-input diagnostic runs or dependence-aware sensitivity methods; Online Appendix G gives the full formulas and protocol.

% S3-SPAN-M10-END
\endgroup

\section{Full result grids}

\begingroup
\def\gridcaption{Full grid backing compact Table 4 of the main text}
\makeatletter
\let\gridoriginalLTmakecaption\LT@makecaption
\renewcommand{\LT@makecaption}[3]{\gridoriginalLTmakecaption{#1}{#2}{\gridcaption}}
\makeatother
\renewcommand{\label}[1]{}
\begingroup
\fontsize{8}{9.6}\selectfont
\setlength{\tabcolsep}{1.2pt}
\let\tableunderscore\_
\renewcommand{\_}{\tableunderscore\allowbreak}
\begin{longtable}{@{}
>{\raggedright\arraybackslash}p{0.06\linewidth}
>{\raggedright\arraybackslash}p{0.23\linewidth}
>{\raggedright\arraybackslash}p{0.10\linewidth}
>{\raggedright\arraybackslash}p{0.07\linewidth}
>{\raggedright\arraybackslash}p{0.19\linewidth}
@{\hspace{0.35em}}
>{\raggedright\arraybackslash}p{0.12\linewidth}
>{\raggedright\arraybackslash}p{0.18\linewidth}
@{}}
\caption{Table 4. Baseline $V^{gross}$: primary-case moderate-scenario block and all crossing cells}
\label{tab:results_vgross}\\
\toprule
country\_id & policy\_variant\_id & scenario\_id & regime\_id & v\_gross & crosses\_v1 & cell\_result\_class \\
\midrule
\endfirsthead
\caption[]{Table 4. Baseline $V^{gross}$: primary-case moderate-scenario block and all crossing cells (continued)}\\
\toprule
country\_id & policy\_variant\_id & scenario\_id & regime\_id & v\_gross & crosses\_v1 & cell\_result\_class \\
\midrule
\endhead
\multicolumn{7}{l}{\textit{Panel A: Peru, moderate scenario, all policy variants and regimes}} \\
\midrule
PER & PEN & mid & r0 & 0.0781898593701847 & False & not\_feasible \\
PER & GMI:GMI\_ideal\_aggregate & mid & r0 & -0.0525921743988333 & False & not\_feasible \\
PER & GMI:GMI\_loaded\_aggregate & mid & r0 & -0.0350614495992222 & False & not\_feasible \\
PER & MUT & mid & r0 & 0.0191207387806886 & False & not\_feasible \\
PER & PBI & mid & r0 & 0.0164818464107766 & False & not\_feasible \\
PER & UBI & mid & r0 & 0.0068349135400345 & False & not\_feasible \\
PER & PEN & mid & r1 & 0.104105996516133 & False & not\_feasible \\
PER & GMI:GMI\_ideal\_aggregate & mid & r1 & -0.0464766449646969 & False & not\_feasible \\
PER & GMI:GMI\_loaded\_aggregate & mid & r1 & -0.0309844299764646 & False & not\_feasible \\
PER & MUT & mid & r1 & 0.0216762789374418 & False & not\_feasible \\
PER & PBI & mid & r1 & 0.018143849607453 & False & not\_feasible \\
PER & UBI & mid & r1 & 0.0074457059524051 & False & not\_feasible \\
PER & PEN & mid & r2 & 0.1232497556668719 & False & not\_feasible \\
PER & GMI:GMI\_ideal\_aggregate & mid & r2 & -0.042648285425488 & False & not\_feasible \\
PER & GMI:GMI\_loaded\_aggregate & mid & r2 & -0.0284321902836587 & False & not\_feasible \\
PER & MUT & mid & r2 & 0.0236746826169089 & False & not\_feasible \\
PER & PBI & mid & r2 & 0.0194827667412041 & False & not\_feasible \\
PER & UBI & mid & r2 & 0.0079453068722719 & False & not\_feasible \\
PER & PEN & mid & r3 & 0.0721660346552839 & False & not\_feasible \\
PER & GMI:GMI\_ideal\_aggregate & mid & r3 & -0.0534236826727717 & False & not\_feasible \\
PER & GMI:GMI\_loaded\_aggregate & mid & r3 & -0.0356157884485144 & False & not\_feasible \\
PER & MUT & mid & r3 & 0.0184319849785464 & False & not\_feasible \\
PER & PBI & mid & r3 & 0.0160003092940026 & False & not\_feasible \\
PER & UBI & mid & r3 & 0.0066514880387719 & False & not\_feasible \\
PER & PEN & mid & r4 & 0.1765945556483395 & False & not\_feasible \\
PER & GMI:GMI\_ideal\_aggregate & mid & r4 & -0.0310475854881461 & False & not\_feasible \\
PER & GMI:GMI\_loaded\_aggregate & mid & r4 & -0.0206983903254307 & False & not\_feasible \\
PER & MUT & mid & r4 & 0.0290934817763354 & False & not\_feasible \\
PER & PBI & mid & r4 & 0.023063134368669 & False & not\_feasible \\
PER & UBI & mid & r4 & 0.0092719140353109 & False & not\_feasible \\
\midrule
\multicolumn{7}{l}{\textit{Panel B: all crossing cells}} \\
\midrule
CHL & GMI:GMI\_ideal\_aggregate & stress & r0 & 3.901937871138183 & True & feasible\_cell \\
CHL & GMI:GMI\_ideal\_aggregate & stress & r1 & 4.239349420616151 & True & feasible\_cell \\
CHL & GMI:GMI\_ideal\_aggregate & stress & r2 & 4.550059741080454 & True & feasible\_cell \\
CHL & GMI:GMI\_ideal\_aggregate & stress & r3 & 4.015460658881374 & True & feasible\_cell \\
CHL & GMI:GMI\_ideal\_aggregate & stress & r4 & 5.799668762204835 & True & feasible\_cell \\
CHL & GMI:GMI\_loaded\_aggregate & stress & r0 & 2.6012919140921213 & True & feasible\_cell \\
CHL & GMI:GMI\_loaded\_aggregate & stress & r1 & 2.826232947077433 & True & feasible\_cell \\
CHL & GMI:GMI\_loaded\_aggregate & stress & r2 & 3.033373160720303 & True & feasible\_cell \\
CHL & GMI:GMI\_loaded\_aggregate & stress & r3 & 2.676973772587582 & True & feasible\_cell \\
CHL & GMI:GMI\_loaded\_aggregate & stress & r4 & 3.86644584146989 & True & feasible\_cell \\
PER & PEN & stress & r0 & 1.0152041145854496 & True & accounting\_feasible\_debt\_failed \\
PER & PEN & stress & r1 & 1.0926140795425037 & True & accounting\_feasible\_debt\_failed \\
PER & PEN & stress & r2 & 1.15780271593806 & True & accounting\_feasible\_debt\_failed \\
PER & PEN & stress & r3 & 1.0499800745562868 & True & accounting\_feasible\_debt\_failed \\
PER & PEN & stress & r4 & 1.4222907379520815 & True & accounting\_feasible\_debt\_failed \\
\bottomrule
\end{longtable}
\endgroup

\endgroup

\begingroup
\def\gridcaption{Full grid backing compact Table 5 of the main text}
\makeatletter
\let\gridoriginalLTmakecaption\LT@makecaption
\renewcommand{\LT@makecaption}[3]{\gridoriginalLTmakecaption{#1}{#2}{\gridcaption}}
\makeatother
\renewcommand{\label}[1]{}
\begin{landscape}
\fontsize{6}{7.2}\selectfont
\setlength{\tabcolsep}{1.5pt}
\setlength{\LTleft}{0pt}
\setlength{\LTright}{0pt}
\let\tableunderscore\_
\renewcommand{\_}{\tableunderscore\allowbreak}
\begin{longtable}{@{}
>{\raggedright\arraybackslash}p{0.03\linewidth}
>{\raggedright\arraybackslash}p{0.13\linewidth}
>{\raggedright\arraybackslash}p{0.045\linewidth}
>{\raggedright\arraybackslash}p{0.035\linewidth}
>{\raggedright\arraybackslash}p{0.09\linewidth}
>{\raggedright\arraybackslash}p{0.09\linewidth}
>{\raggedright\arraybackslash}p{0.09\linewidth}
>{\raggedright\arraybackslash}p{0.08\linewidth}
>{\raggedright\arraybackslash}p{0.17\linewidth}
>{\raggedright\arraybackslash}p{0.19\linewidth}
@{}}
\caption{Table 5. Threshold inversion: headline cells}\label{tab:results_inversions}\\
\toprule
country\_id & policy\_variant\_id & scenario\_id & regime\_id & requirement\_basis & g\_ai\_required & mfc\_required\_gross & q\_prod\_required & historical\_total\_revenue\_class & flags \\
\midrule
\endfirsthead
\caption[]{Table 5. Threshold inversion: headline cells (continued)}\\
\toprule
country\_id & policy\_variant\_id & scenario\_id & regime\_id & requirement\_basis & g\_ai\_required & mfc\_required\_gross & q\_prod\_required & historical\_total\_revenue\_class & flags \\
\midrule
\endhead
CHL & GMI:GMI\_ideal\_aggregate & mid & r0 & baseline & 0.058961369896838 & 0.1703917298779569 &  & fiscally\_ordinary & frontier\_exceeds\_one;ai\_rent\_capture\_impossible \\
CHL & GMI:GMI\_ideal\_aggregate & mid & r0 & debt\_consistent & 0.058961369896838 & 0.1703917298779569 &  & fiscally\_ordinary & frontier\_exceeds\_one;ai\_rent\_capture\_impossible \\
CHL & GMI:GMI\_ideal\_aggregate & mid & r1 & baseline & 0.055747862766219 & 0.1687742033362891 &  & fiscally\_ordinary & frontier\_exceeds\_one;ai\_rent\_capture\_impossible \\
CHL & GMI:GMI\_ideal\_aggregate & mid & r1 & debt\_consistent & 0.055747862766219 & 0.1687742033362891 &  & fiscally\_ordinary & frontier\_exceeds\_one;ai\_rent\_capture\_impossible \\
CHL & GMI:GMI\_ideal\_aggregate & mid & r2 & baseline & 0.0532949307271395 & 0.1687742033362891 &  & fiscally\_ordinary & frontier\_exceeds\_one;ai\_rent\_capture\_impossible \\
CHL & GMI:GMI\_ideal\_aggregate & mid & r2 & debt\_consistent & 0.0532949307271395 & 0.1687742033362891 &  & fiscally\_ordinary & frontier\_exceeds\_one;ai\_rent\_capture\_impossible \\
CHL & GMI:GMI\_ideal\_aggregate & mid & r3 & baseline & 0.0603464096353531 & 0.1697447192612898 &  & fiscally\_ordinary & frontier\_exceeds\_one;ai\_rent\_capture\_impossible \\
CHL & GMI:GMI\_ideal\_aggregate & mid & r3 & debt\_consistent & 0.0603464096353531 & 0.1697447192612898 &  & fiscally\_ordinary & frontier\_exceeds\_one;ai\_rent\_capture\_impossible \\
CHL & GMI:GMI\_ideal\_aggregate & mid & r4 & baseline & 0.0474832853486911 & 0.1671566767946213 &  & fiscally\_ordinary & frontier\_exceeds\_one;ai\_rent\_capture\_impossible \\
CHL & GMI:GMI\_ideal\_aggregate & mid & r4 & debt\_consistent & 0.0474832853486911 & 0.1671566767946213 &  & fiscally\_ordinary & frontier\_exceeds\_one;ai\_rent\_capture\_impossible \\
CHL & GMI:GMI\_ideal\_aggregate & stress & r0 & baseline & 0.0635180536168128 & 0.0466217878394694 & 0.07573616388301 & fiscally\_ordinary & requirement\_already\_covered \\
CHL & GMI:GMI\_ideal\_aggregate & stress & r0 & debt\_consistent & 0.0635180536168128 & 0.0466217878394694 & 0.07573616388301 & fiscally\_ordinary & requirement\_already\_covered \\
CHL & GMI:GMI\_ideal\_aggregate & stress & r1 & baseline & 0.0600690925713408 & 0.0461792077957997 & 0.071729694873811 & fiscally\_ordinary & requirement\_already\_covered \\
CHL & GMI:GMI\_ideal\_aggregate & stress & r1 & debt\_consistent & 0.0600690925713408 & 0.0461792077957997 & 0.071729694873811 & fiscally\_ordinary & requirement\_already\_covered \\
CHL & GMI:GMI\_ideal\_aggregate & stress & r2 & baseline & 0.0574083189193934 & 0.0461792077957997 & 0.0686307955270862 & fiscally\_ordinary & requirement\_already\_covered \\
CHL & GMI:GMI\_ideal\_aggregate & stress & r2 & debt\_consistent & 0.0574083189193934 & 0.0461792077957997 & 0.0686307955270862 & fiscally\_ordinary & requirement\_already\_covered \\
CHL & GMI:GMI\_ideal\_aggregate & stress & r3 & baseline & 0.0623048192908522 & 0.0464447558220015 & 0.0743281510510759 & fiscally\_ordinary & requirement\_already\_covered \\
CHL & GMI:GMI\_ideal\_aggregate & stress & r3 & debt\_consistent & 0.0623048192908522 & 0.0464447558220015 & 0.0743281510510759 & fiscally\_ordinary & requirement\_already\_covered \\
CHL & GMI:GMI\_ideal\_aggregate & stress & r4 & baseline & 0.04846101111926 & 0.0457366277521299 & 0.058158483206631 & fiscally\_ordinary & requirement\_already\_covered \\
CHL & GMI:GMI\_ideal\_aggregate & stress & r4 & debt\_consistent & 0.04846101111926 & 0.0457366277521299 & 0.058158483206631 & fiscally\_ordinary & requirement\_already\_covered \\
PER & PEN & mid & r0 & baseline & 0.0930780545727182 & 0.3049656645066527 &  & fiscally\_demanding & frontier\_exceeds\_one;ai\_rent\_capture\_impossible \\
PER & PEN & mid & r0 & debt\_consistent & 0.1769546915535554 & 0.5797833371671834 &  & extreme\_or\_outside\_historical\_support & frontier\_exceeds\_one;ai\_rent\_capture\_impossible \\
PER & PEN & mid & r1 & baseline & 0.0879595921251269 & 0.3040376146025564 &  & fiscally\_demanding & frontier\_exceeds\_one;ai\_rent\_capture\_impossible \\
PER & PEN & mid & r1 & debt\_consistent & 0.1674657099046527 & 0.5788552872630872 &  & extreme\_or\_outside\_historical\_support & frontier\_exceeds\_one;ai\_rent\_capture\_impossible \\
PER & PEN & mid & r2 & baseline & 0.0840785922521413 & 0.3040376146025564 &  & fiscally\_demanding & frontier\_exceeds\_one;ai\_rent\_capture\_impossible \\
PER & PEN & mid & r2 & debt\_consistent & 0.1600766988466565 & 0.5788552872630872 &  & extreme\_or\_outside\_historical\_support & frontier\_exceeds\_one;ai\_rent\_capture\_impossible \\
PER & PEN & mid & r3 & baseline & 0.094833365540356 & 0.3045944445450142 &  & fiscally\_demanding & frontier\_exceeds\_one;ai\_rent\_capture\_impossible \\
PER & PEN & mid & r3 & debt\_consistent & 0.1803959398916242 & 0.579412117205545 &  & extreme\_or\_outside\_historical\_support & frontier\_exceeds\_one;ai\_rent\_capture\_impossible \\
PER & PEN & mid & r4 & baseline & 0.0753100315197287 & 0.3031095646984601 &  & fiscally\_demanding & frontier\_exceeds\_one;ai\_rent\_capture\_impossible \\
PER & PEN & mid & r4 & debt\_consistent & 0.1435907128332082 & 0.5779272373589909 &  & extreme\_or\_outside\_historical\_support & frontier\_exceeds\_one;ai\_rent\_capture\_impossible \\
PER & PEN & stress & r0 & baseline & 0.0999721807025207 & 0.0839172321326687 & 0.1173826778266382 & fiscally\_ordinary &  \\
PER & PEN & stress & r0 & debt\_consistent & 0.1900614111603496 & 0.1595386581319568 &  & fiscally\_ordinary & frontier\_exceeds\_one \\
PER & PEN & stress & r1 & baseline & 0.0944755888991784 & 0.0836618611571895 & 0.1111834472512639 & fiscally\_ordinary &  \\
PER & PEN & stress & r1 & debt\_consistent & 0.1798714748603459 & 0.1592832871564775 &  & fiscally\_ordinary & frontier\_exceeds\_one \\
PER & PEN & stress & r2 & baseline & 0.090281494348363 & 0.0836618611571895 & 0.1064343388387812 & fiscally\_ordinary & requirement\_already\_covered \\
PER & PEN & stress & r2 & debt\_consistent & 0.1718863648298174 & 0.1592832871564775 &  & fiscally\_ordinary & frontier\_exceeds\_one \\
PER & PEN & stress & r3 & baseline & 0.0974293898017232 & 0.083815083742477 & 0.1145183113777666 & fiscally\_ordinary &  \\
PER & PEN & stress & r3 & debt\_consistent & 0.185334204329909 & 0.159436509741765 &  & fiscally\_ordinary & frontier\_exceeds\_one \\
PER & PEN & stress & r4 & baseline & 0.0764586709452878 & 0.0834064901817102 & 0.0906650576363045 & fiscally\_ordinary & requirement\_already\_covered \\
PER & PEN & stress & r4 & debt\_consistent & 0.1457807790246047 & 0.1590279161809982 & 0.1679894898554811 & fiscally\_ordinary &  \\
\bottomrule
\end{longtable}
\end{landscape}

\endgroup

\begingroup
\def\gridcaption{Full grid backing compact Table 6 of the main text}
\makeatletter
\let\gridoriginalLTmakecaption\LT@makecaption
\renewcommand{\LT@makecaption}[3]{\gridoriginalLTmakecaption{#1}{#2}{\gridcaption}}
\makeatother
\renewcommand{\label}[1]{}
\input{tables/table6_supp_full}
\endgroup

\begingroup
\def\gridcaption{Full grid backing compact Table 7 of the main text}
\makeatletter
\let\gridoriginalLTmakecaption\LT@makecaption
\renewcommand{\LT@makecaption}[3]{\gridoriginalLTmakecaption{#1}{#2}{\gridcaption}}
\makeatother
\renewcommand{\label}[1]{}
\input{tables/table7_supp_full}
\endgroup

\begin{thebibliography}{99}
\bibitem[Acemoglu(2024)]{acemoglu2024simple}
Acemoglu, D. (2024)
\newblock The simple macroeconomics of AI.
\newblock \emph{NBER Working Paper No. 32487}. DOI: 10.3386/w32487.

\bibitem[Acemoglu and Restrepo(2018)]{acemoglu2018race}
Acemoglu, D. and Restrepo, P. (2018)
\newblock The race between man and machine: Implications of technology for growth, factor shares, and employment.
\newblock \emph{American Economic Review} 108(6): 1488--1542. DOI: 10.1257/aer.20160696.

\bibitem[Acemoglu and Restrepo(2019)]{acemoglu2019automation}
Acemoglu, D. and Restrepo, P. (2019)
\newblock Automation and new tasks: How technology displaces and reinstates labor.
\newblock \emph{Journal of Economic Perspectives} 33(2): 3--30. DOI: 10.1257/jep.33.2.3.


\bibitem[Athey and Imbens(2017)]{athey2017state}
Athey, S. and Imbens, G.W. (2017)
\newblock The state of applied econometrics: Causality and policy evaluation.
\newblock \emph{Journal of Economic Perspectives} 31(2): 3--32. DOI: 10.1257/jep.31.2.3.

\bibitem[Bass(1969)]{bass1969newproduct}
Bass, F.M. (1969)
\newblock A new product growth for model consumer durables.
\newblock \emph{Management Science} 15(5): 215--227. DOI: 10.1287/mnsc.15.5.215.

\bibitem[Bertrand, Duflo and Mullainathan(2004)]{bertrand2004trust}
Bertrand, M., Duflo, E., and Mullainathan, S. (2004)
\newblock How much should we trust differences-in-differences estimates?
\newblock \emph{Quarterly Journal of Economics} 119(1): 249--275. DOI: 10.1162/003355304772839588.

\bibitem[Brollo et al.(2024)]{brollo2024fiscal}
Brollo, F., Dabla-Norris, E., de Mooij, R., Garcia-Macia, D., Hanappi, T., Liu, L., and Nguyen, A.D.M. (2024)
\newblock Broadening the gains from generative AI: The role of fiscal policies.
\newblock \emph{IMF Staff Discussion Note 2024/002}. DOI: 10.5089/9798400277177.006.

\bibitem[Brynjolfsson, Li and Raymond(2023)]{brynjolfsson2023generative}
Brynjolfsson, E., Li, D., and Raymond, L.R. (2023)
\newblock Generative AI at work.
\newblock \emph{NBER Working Paper No. 31161}. DOI: 10.3386/w31161.

\bibitem[Cazzaniga et al.(2024)]{cazzaniga2024genai}
Cazzaniga, M., Jaumotte, F., Li, L., Melina, G., Panton, A.J., Pizzinelli, C., Rockall, E.J., and Tavares, M.M. (2024)
\newblock Gen-AI: Artificial intelligence and the future of work.
\newblock \emph{IMF Staff Discussion Note 2024/001}. DOI: 10.5089/9798400262548.006.



\bibitem[Filippucci et al.(2024)]{filippucci2024miracle}
Filippucci, F., Gal, P., and Schief, M. (2024)
\newblock Miracle or myth? Assessing the macroeconomic productivity gains from artificial intelligence.
\newblock \emph{OECD Artificial Intelligence Papers No. 29}. DOI: 10.1787/b524a072-en.


\bibitem[Filippucci et al.(2025)]{oecd2025g7}
Filippucci, F., Gal, P., Laengle, K., and Schief, M. (2025)
\newblock Macroeconomic productivity gains from artificial intelligence in G7 economies.
\newblock \emph{OECD Artificial Intelligence Papers No. 41}. DOI: 10.1787/a5319ab5-en.


\bibitem[Gentilini et al.(2020)]{gentilini2020ubi}
Gentilini, U., Grosh, M., Rigolini, J., and Yemtsov, R. (2020)
\newblock \emph{Exploring Universal Basic Income: A Guide to Navigating Concepts, Evidence, and Practices}.
\newblock World Bank. DOI: 10.1596/978-1-4648-1458-7.


\bibitem[ILOSTAT(2026)]{ilostat2026informal}
ILOSTAT (2026)
\newblock Informal employment rate.
\newblock International Labour Organization statistical snapshot and database, accessed June 2026.

\bibitem[International Monetary Fund(2024)]{imf2024aipi}
International Monetary Fund (2024)
\newblock AI Preparedness Index (AIPI).
\newblock IMF DataMapper dataset and documentation.

\bibitem[Lempert, Popper and Bankes(2003)]{lempert2003shaping}
Lempert, R.J., Popper, S.W., and Bankes, S.C. (2003)
\newblock \emph{Shaping the Next One Hundred Years: New Methods for Quantitative, Long-Term Policy Analysis}.
\newblock RAND Corporation, MR-1626-RPC.



\bibitem[Lipsitch, Tchetgen Tchetgen and Cohen(2010)]{lipsitch2010negative}
Lipsitch, M., Tchetgen Tchetgen, E., and Cohen, T. (2010)
\newblock Negative controls: A tool for detecting confounding and bias in observational studies.
\newblock \emph{Epidemiology} 21(3): 383--388. DOI: 10.1097/EDE.0b013e3181d61eeb.

\bibitem[McKay, Beckman and Conover(1979)]{mckay1979lhs}
McKay, M.D., Beckman, R.J., and Conover, W.J. (1979)
\newblock A comparison of three methods for selecting values of input variables in the analysis of output from a computer code.
\newblock \emph{Technometrics} 21(2): 239--245. DOI: 10.2307/1268522.

\bibitem[Medina and Schneider(2018)]{medina2018shadow}
Medina, L. and Schneider, F. (2018)
\newblock Shadow economies around the world: What did we learn over the last 20 years?
\newblock \emph{IMF Working Paper 18/17}. DOI: 10.5089/9781484338636.001.

\bibitem[Morris(1991)]{morris1991factorial}
Morris, M.D. (1991)
\newblock Factorial sampling plans for preliminary computational experiments.
\newblock \emph{Technometrics} 33(2): 161--174. DOI: 10.1080/00401706.1991.10484804.

\bibitem[Noy and Zhang(2023)]{noy2023experimental}
Noy, S. and Zhang, W. (2023)
\newblock Experimental evidence on the productivity effects of generative artificial intelligence.
\newblock \emph{Science} 381(6654): 187--192. DOI: 10.1126/science.adh2586.

\bibitem[OECD et al.(2025)]{oecd2025revenue}
OECD, ECLAC, CIAT, and IDB (2025)
\newblock \emph{Revenue Statistics in Latin America and the Caribbean 2025}.
\newblock OECD Publishing. DOI: 10.1787/7594fbdd-en.

\bibitem[RAND Corporation(2013)]{rand2013rdm}
RAND Corporation (2013)
\newblock Making good decisions without predictions: Robust Decision Making for planning under deep uncertainty.
\newblock RAND Research Highlights, RB-9701.


\bibitem[Rogers(2003)]{rogers2003diffusion}
Rogers, E.M. (2003)
\newblock \emph{Diffusion of Innovations}, 5th ed.
\newblock Free Press.



\bibitem[Saltelli et al.(2008)]{saltelli2008global}
Saltelli, A., Ratto, M., Andres, T., Campolongo, F., Cariboni, J., Gatelli, D., Saisana, M., and Tarantola, S. (2008)
\newblock \emph{Global Sensitivity Analysis: The Primer}.
\newblock Wiley.

\bibitem[Sobol(2001)]{sobol2001global}
Sobol, I.M. (2001)
\newblock Global sensitivity indices for nonlinear mathematical models and their Monte Carlo estimates.
\newblock \emph{Mathematics and Computers in Simulation} 55(1--3): 271--280. DOI: 10.1016/S0378-4754(00)00270-6.

\bibitem[UNU-WIDER(2025)]{unu2025grd}
UNU-WIDER (2025)
\newblock Government Revenue Dataset, Version 2025.
\newblock DOI: 10.35188/UNU-WIDER/GRD-2025.

\bibitem[World Bank(2021)]{worldbank2021informality}
World Bank (2021)
\newblock \emph{The Long Shadow of Informality: Challenges and Policies}.
\newblock Washington, DC: World Bank. DOI: 10.1596/978-1-4648-1753-3.


\bibitem[World Bank(2026a)]{worldbank2026aspire}
World Bank (2026a)
\newblock ASPIRE: The Atlas of Social Protection Indicators of Resilience and Equity.
\newblock Social protection and labor indicators, accessed June 2026.

\bibitem[World Bank(2026b)]{worldbank2026pip}
World Bank (2026b)
\newblock Poverty and Inequality Platform.
\newblock Poverty, inequality, and poverty-line data, accessed June 2026.

\bibitem[World Bank and ILO(2024)]{worldbankilo2024lacgenai}
World Bank and International Labour Organization (2024)
\newblock Generative AI and jobs in Latin America and the Caribbean: Is the digital divide a buffer or bottleneck?
\newblock World Bank Group and International Labour Organization.

\end{thebibliography}

\end{document}
